\clearpage

\section{\MakeUppercase{РОЗРОБЛЕННЯ НАВЧАЛЬНО-МЕТОДИЧНИХ МАТЕРІАЛІВ ДО ДИФЕРЕНЦІЙОВАНОГО НАВЧАННЯ ВИБІРКОВОГО МОДУЛЯ ``ІНФОРМАЦІЙНА БЕЗПЕКА'' УЧНІВ ЛІЦЕЇВ}}

\subsection{Обґрунтування змістовного наповнення вибіркового модуля ``Інформаційна безпека'' учнів ліцею}




\subsubsection{Порівняльний аналіз програми навчання учнів інформаційної безпеки}
%В умовах зростання обсягів даних, повсюдної інтеграції цифрових технологій та становлення інформаційного суспільства інформація набуває статусу ключового стратегічного ресурсу та економічного активу. Відповідно, забезпечення її захищеності та життєздатності, відоме як інформаційна безпека,  стає імперативом для стійкого функціонування будь-якого суб'єкта~-- від окремої особистості до глобальних корпорацій та держав. Вивчення основ інформаційної безпеки є фундаментальним, системним кроком у формуванні міждисциплінарних компетентностей, необхідних для ефективної ідентифікації, оцінки та нейтралізації сучасних загроз у динамічному кіберпросторі. 

%Необхідність вивчення основ інформаційної безпеки базується на трьох основних групах факторів: актуальність у сучасному світі, захист особистих та корпоративних інтересів та професійна затребуваність.

%Актуальність у сучасному світі. У сучасному світі, що динамічно цифровізується, інформація є одним із найцінніших ресурсів, активом як для приватних осіб, так і для організацій та держав, оскільки вона зберігається, обробляється та передається переважно в цифровому вигляді. Із розвитком технологій відбувається постійне зростання кіберзагроз, збільшується кількість, складність та інтенсивність кіберзлочинів (віруси-вимагачі, фішинг, хакерські атаки, шпигунство), і тому інформаційна безпека є критично необхідною для запобігання цим загрозам та мінімізації їхніх наслідків. Окрім того, державний та національний рівень безпеки неможливий без інформаційної безпеки, що включає захист критичної інфраструктури (енергетика, фінанси, транспорт), державних реєстрів та протидію дезінформації та інформаційним війнам. Нарешті, законодавче регулювання (включаючи українські закони та міжнародні стандарти, як-от європейський GDPR) вимагає від організацій забезпечення належного рівня інформаційної безпеки.

%Захист особистих та корпоративних інтересів. Фундаментальне значення інформаційної безпеки закріплене в міжнародних стандартах (зокрема, ISO/IEC 27001) та національному законодавстві, які визначають її через забезпечення трьох ключових компонентів: 

%\textit{конфіденційність (Confidentiality)} полягає у захисті інформації від несанкціонованого доступу, що стосується особистих даних (паролі, фінанси, медичні записи) та корпоративної таємниці (бізнес-плани, клієнтські бази). Забезпечення того, що доступ до інформації можуть отримати виключно авторизовані суб'єкти. Деталізація охоплює методи контролю доступу (зокрема, Role-Based Access Control), механізми шифрування даних (симетричні та асиметричні алгоритми) та управління ідентифікацією (ідентифікація та аутентифікація). Це критично важливо для захисту персональних даних, комерційної таємниці та державної інформації з обмеженим доступом. Вивчення таких питань формує навички захисту приватності та комерційної таємниці. 

%\textit{цілісність (Integrity) }полягає у гарантуванні точності, повноти та незмінності інформації протягом усього її життєвого циклу. Цілісність забезпечує достовірність та незмінність інформації, що є критично важливим для фінансових транзакцій, юридичних документів та державних даних. Вивчення цілісності передбачає освоєння механізмів контрольних сум (хешування), цифрових підписів та протоколів транзакцій, які дозволяють виявляти та запобігати неправомірному чи випадковому модифікуванню даних. Вивчення цих питань допомагає зрозуміти, як виявляти та запобігати несанкціонованим змінам.

%\textit{доступність (Availability)} гарантує, що авторизовані користувачі матимуть безперешкодний доступ до інформації та систем у разі необхідності.   Цей аспект вимагає вивчення принципів відмовостійкості, резервного копіювання, кластеризації та архітектури систем, стійких до кібератак, зокрема атак типу відмови в обслуговуванні (Distributed Denial of Service, DDoS), планування стійкості систем та резервне копіювання для запобігання збоям. Вивчення і практичне опанування цих питань забезпечать своєчасний, надійний та безперешкодний доступ авторизованих користувачів та систем до інформаційних ресурсів.

%Професійна затребуваність. Знання з інформаційної безпеки необхідні кожному, оскільки особиста кібергігієна вимагає від користувача мережі знання основ безпечної поведінки (безпечні паролі, розпізнавання фішингу, оновлення ПЗ) для захисту себе та своєї родини. Водночас, галузь інформаційної (кібер-) безпеки є однією з найбільш швидкозростаючих та високооплачуваних у світі, що створює широкі кар'єрні можливості, адже попит на кваліфікованих фахівців (тестувальників на проникнення, аналітиків загроз, аудиторів безпеки) значно перевищує пропозицію. Крім того, вивчення теми сприяє формуванню культури безпеки, оскільки для будь-якої організації важливо, щоб усі співробітники розуміли свою роль у захисті інформації, формуючи відповідальне ставлення до даних.

%Отже, вивчення інформаційної безпеки  є життєво необхідною навичкою для існування в сучасному інформаційному суспільстві, що дає змогу ефективно захищатися від загроз, відповідати вимогам законодавства та забезпечувати стабільність як особистого життя, так і діяльності організацій.

Інформаційна безпека не є ізольованим модулем, а є наскрізною вимогою курсу інформатики, яка формується через декілька міжпредметних компетентностей, що відображає сучасні педагогічні підходи до формування відповідального цифрового громадянства. Проаналізуємо Державний стандарт базової середньої освіти \cite{KMU_898_2020}, зокрема  інформатичну освітню галузь.  У таблиці~\ref{tab:vimogi_bezpeki} подані вимоги до кібербезпеки, цифрової етики та захисту даних, які інтегровані в структуру ключових компетентностей.

\begin{longtable}{|m{0.05\textwidth}|p{0.25\textwidth}|p{0.6\textwidth}|}

\caption{Вимоги до учнів щодо формування навичок безпечного та відповідального поводження у цифровому середовищі згідно з Державним стандартом базової середньої освіти \cite{KMU_898_2020}}
\label{tab:vimogi_bezpeki} 
% Налаштовуємо вигляд маркованих списків всередині таблиці
\setlist[itemize]{nosep, leftmargin=*, label=\textbullet}
% m{...} - вертикальне вирівнювання по центру, p{...} - по верху
\\
\hline
\textbf{№} & \textbf{\makecell{Ключова\\ компетентність}} & \hfil \textbf{Уміння та ставлення} \\ 
\hline
\endfirsthead % Заголовок для першої сторінки

\hline
\endhead % Заголовок для наступних сторінок

% 1. Екологічна грамотність
1 & \textbf{Екологічна грамотність і здорове життя} & 
\begin{itemize}
    \item Уміння: дотримуватися правил безпечного поводження з комп’ютерними пристроями (апаратними, мережевими та програмними засобами).
    \item Уміння: обирати безпечні паролі та інші методи автентифікації.
\end{itemize} \\ 
\hline

% 2. Громадянська та соціальна
2 & \textbf{Громадянська та соціальна компетентності} & 
\begin{itemize}
    \item Уміння: пояснювати співвідношення між правилами, потребами і правом та законами в галузі цифрових технологій.
    \item Уміння: дотримуватися у власній інформаційній діяльності законів щодо захисту даних, інтелектуальної і приватної власності, пояснювати відповідальність за їх порушення.
    \item Уміння: наводити приклади наслідків порушення прав інтелектуальної власності, розрізняти різновиди і серйозність порушень.
    \item Ставлення: повага до інтелектуальної власності та нетерпимість до порушень правових норм.
\end{itemize} \\ 
\hline

% 3. Інформаційно-комунікаційна (ІКТ)
%3 & \textbf{Інформаційно-комунікаційна компетентність (ІКТ)} & 
%\begin{itemize}
%    \item Уміння: Розрізняти особисту, приватну і конфіденційну інформацію.
%    \item Уміння: Використовувати цифрові інструменти для захисту даних та приватності (наприклад, шифрування, двофакторна автентифікація).
%    \item Уміння: Розпізнавати та уникати фішингових посилань, спаму та іншого небезпечного контенту.
%    \item Уміння: Дотримуватися етикету (нетикету) під час роботи в Інтернеті та обміну повідомленнями.
%    \item Уміння: Знати, як діяти у разі виявлення кіберзагрози або порушення безпеки (наприклад, зараження вірусом, витік даних).
%\end{itemize} \\ 
%\hline

% 4. Навчання впродовж життя
%4 & \textbf{Навчання впродовж життя (як частина ІКТ)} & 
%\begin{itemize}
%    \item Уміння: Оцінювати ризики для приватності та безпеки, пов'язані з використанням цифрових технологій, та обирати параметри безпеки відповідно до особистих потреб.
%    \item Уміння: Використовувати антивірусні та інші захисні програми.
%    \item Уміння: Пояснювати наслідки використання неліцензійного програмного забезпечення.
%\end{itemize} \\ 
%\hline

% 5. Уміння вчитися
%5 & \textbf{Уміння вчитися} & 
%\begin{itemize}
%    \item Уміння: Розрізняти надійні та ненадійні джерела інформації в Інтернеті (що є основою інформаційної безпеки).
%\end{itemize} \\ 
%\hline

\end{longtable}

Як видно з таблиці компетентність ``Екологічна грамотність і здорове життя'' містить важливий блок вимог, що стосуються фізичної та технічної безпеки. Зокрема, стандарт вимагає від здобувачів освіти дотримання правил безпечного поводження з комп’ютерними пристроями (апаратними, мережевими та програмними засобами). Критично важливим є вміння обирати безпечні паролі та інші методи автентифікації, що є базовим елементом захисту персональних облікових записів.

Соціальний вимір кібербезпеки є фундаментальною складовою компетентності ``Громадянська та соціальна компетентності''. Основний аспект цієї компетентності зміщується з ``як захистити себе'' на ``як діяти відповідально та законно''. Учні повинні не просто знати, а й пояснювати співвідношення між правилами, потребами і правом та законами в галузі цифрових технологій. Центральною вимогою є дотримання законів щодо захисту даних, інтелектуальної і приватної власності у власній інформаційній діяльності. Стандарт вимагає від учнів здатності пояснювати відповідальність за їх порушення, включно з умінням наводити приклади наслідків порушення прав інтелектуальної власності та розрізняти серйозність цих порушень \cite{KMU_898_2020}. Це формує повагу до інтелектуальної власності та нетерпимість до порушень правових норм.

Ядро практичного кіберзахисту зосереджено в компетентності ``Інформаційно-комунікаційна компетентність (ІКT)''. Тут формуються навички активного управління власним цифровим слідом та захистом. Це охоплює здатність розрізняти особисту, приватну і конфіденційну інформацію, що є необхідною передумовою для усвідомленого захисту. Технічний захист передбачає використання цифрових інструментів для захисту даних, включаючи освоєння шифрування та двофакторної автентифікації. Не менш важливим є превентивний захист: учні мають вміти розпізнавати та уникати фішингових посилань, спаму та іншого небезпечного контенту. Нарешті, Стандарт вимагає знання протоколів поведінки: дотримання етикету (нетикету) в мережевій комунікації та знання алгоритму дій у разі виявлення кіберзагрози чи витоку даних \cite{KMU_898_2020}.

Вимоги, що стосуються безпеки, пронизують і наскрізні компетентності. У межах ``Навчання впродовж життя'' учні мають набути навичок динамічної безпеки: оцінювати ризики для приватності та безпеки при освоєнні нових технологій \cite{KMU_898_2020}. Це також включає практичне вміння використовувати антивірусні та інші захисні програми та розуміння наслідків використання неліцензійного програмного забезпечення. 

Компетентність ``Уміння вчитися'' формує основу для інформаційної безпеки, вимагаючи від учнів розрізняти надійні та ненадійні джерела інформації в Інтернеті \cite{KMU_898_2020}. Це є критично важливим для протидії дезінформації та соціальній інженерії.

Державний стандарт України визначає компетентнісний потенціал Інформатичної освітньої галузі, виділяючи ключові вміння та ставлення, які мають бути сформовані у здобувачів освіти. Проте, стандарт не є переліком тем; він є нормативним каркасом і філософською основою. Завдання навчальної програми полягає у трансформації цих загальних вимог у конкретний, структурований і послідовний навчальний зміст. 

З інформаційною безпекою учні знайомляться ще у базовій школі.  Проаналізуємо більш детально модельні навчальні програми ``Інформатика. 5-6 клас'' \cite{inform_5_6_rivkind}  та ``Інформатика. 7-9 клас'' \cite{inform_7_9_rivkind} для закладів загальної середньої освіти (авторів  Ривкінд Й. Я., Лисенко Т. І., Чернікова~Л. А., Шакотько В. В.), які побудовані на концепції Нової української школи \cite{nush_mon_compressed}. 

Модельна навчальна програма ``Інформатика. 5-6 класи'' \cite{inform_5_6_rivkind} вже на етапі пояснювальної записки визначає безпеку як ключовий елемент, де у меті навчання зазначається: ``
\textit{...розвиток особистості учня, здатного... безпечно та відповідально діяти в інформаційному суспільстві.}''

Модельна навчальна програма ``Інформатика'' для 5–6 класів формує базову цифрову культуру, концентруючись на етичних, правових та гігієнічних аспектах безпеки. На цьому етапі навчання спрямоване на виховання відповідального користувача, здатного безпечно взаємодіяти з початковими цифровими інструментами. Мережевий етикет є наріжним каменем програми 5–6 класів. Учні не просто знайомляться з нетікетом, а засвоюють правила культурної та поважної взаємодії в онлайн-середовищах. Це включає розуміння тону спілкування, уникнення конфліктів та формування позитивної цифрової репутації. Ця складова є фундаментальною для запобігання соціальним загрозам (наприклад, кібербулінгу). З моменту створення перших інформаційних продуктів (презентацій, текстових документів), від учнів вимагається розуміння інтелектуальної власності. Вони мають розрізняти різні типи дозволів на використання контенту та обов'язково зазначати джерела використаних матеріалів. Це не лише правова вимога, а й перший практичний крок до забезпечення цілісності даних (використання перевіреної інформації) та академічної чесності. Під вимогою ``Дотримується правил кібербезпеки'' мається на увазі опанування набору критично важливих поведінкових правил. Це стосується використання надійних, унікальних паролів, обережного ставлення до невідомих електронних листів і посилань (первинна профілактика фішингу), а також розуміння, коли та яку особисту інформацію можна розголошувати. Хоча навчальна програма не вивчає резервне копіювання, тема ``Інформаційні процеси та системи'' формує розуміння цінності даних через їх організацію, збереження та управління файловою структурою. Це є непрямою, але необхідною передумовою для подальшого вивчення системного захисту. Отже, у 5-6 класах за модельною навчальною програмою  закладається  фундамент комп’ютерної безпеки концентруючись на формуванні відповідальної та етичної цифрової поведінки, але не торкається технічної складової захисту.

Модельна навчальна програма ``Інформатика'' для 7–9 класів \cite{inform_7_9_rivkind} значно поглиблює знання, які були закладені раніше, особливо в частині захисту персональних даних та приватності. Крім того відображає зростання соціальної та інформаційної складності, посилюючи фокус на критичному мисленні та відповідальності в мережевих середовищах. У темі ``Інформаційні процеси та системи. Інформація'' (7 клас), учні переходять від простого пошуку інформації до її оцінювання. Вони вчаться виявляти, аналізувати та спростовувати фейки, маніпулятивну та недостовірну інформацію. Це є життєво важливою оборонною стратегією інформаційної безпеки, спрямованою на захист від інформаційних атак. При вивченні комунікації у 7 класі, учні вже вчаться захищати паролі та персональні дані. Тема ``Комп’ютерні мережі. Мережеві технології'' (8 клас) виводить загальні правила кібербезпеки (з 5-6 класів) на більш високий рівень. Учні вчаться аналізувати небезпеки під час використання конкретних мережевих сервісів та усвідомлювати наслідки своїх дій у більших мережевих спільнотах. В рамках опрацювання інформації, виникає потреба захищати персональні дані від несанкціонованого доступу та розуміти основи шифрування. Це включає непряме розуміння важливості аутентифікації та захисту своїх облікових записів. Теми авторського права та етичних норм під час спілкування (8 клас) повторюються, підтверджуючи їхню незмінну важливість у будь-якому віці, особливо в контексті збільшення обсягу створюваного та споживаного контенту. У 9 класі, при вивченні ``Інформаційних систем'', учні активно використовують хмарні сервіси для колективної роботи. Це вимагає від них дотримання правил безпеки в спільних та віддалених середовищах.  На рівні проєктної діяльності (9 клас) поглиблюється ціннісна складова щодо дотримання правил безпеки та запобігання конфліктам. Таким чином,  у 7-9 класах програма виводить поняття безпеки за межі персонального пристрою, підводячи до концепцій мережевої взаємодії та захисту даних, що зберігаються поза контролем користувача.

У навчальній програмі з інформатики для учнів 10-11 класів \cite{inform_10_11_standart} закладаються основи, які орієнтуються на поведінкові аспекти та широку інформаційну грамотність. Базовий модуль переважно зосереджений на поведінковій безпеці та концепції цифрового громадянства, навчаючи учнів загальним правилам безпечної роботи та етиці в мережі Інтернет. Його основна мета полягає у формуванні знань про те, ``як безпечно користуватися'' інформаційними технологіями, а також охоплює широкий спектр тем: від штучного інтелекту та Інтернету речей до фінансових розрахунків і \text{СУБД}. Безпека тут є лише одним із елементів, що підтримує загальну компетентність. У питаннях загроз програма базового модулю обмежується лише загальними ``проблемами та загрозами в Інтернеті'', що є лише поверхневим знайомством. Учні дізнаються про ризики, але не про механізми їхньої реалізації. Щодо засобів захисту, то базовий модуль програми вимагає від учнів лише дотримання правил безпечної поведінки.

Питання основних понять інформаційної безпеки та загроз представлені на рівні знаннєвої складової цієї лінії. Тут вводиться необхідна термінологія, яка дозволяє учням визначити загальну проблематику інформаційної безпеки і сформувати уявлення про роль інформаційної системи та потенційні небезпеки. Учні дізнаються про основи інформаційної безпеки та проблеми інформаційної безпеки, що є відправною точкою для поглибленого вивчення.

Критично важливим є питання дотримання правил безпечної роботи. Програма базового модулю, у своїй діяльнісній складовій, чітко вимагає, щоб учень дотримувався правил безпечної поведінки в Інтернеті та використовував технології цифрового громадянства. Це забезпечує акцент на етиці, особистій відповідальності, мережевому етикету та усвідомленій поведінці в мережі. Це є непрямим охопленням захисту від шкідливого програмного забезпечення, оскільки безпечна поведінка є першим і найважливішим кроком у запобіганні завантаженню вірусів, хоча й без вивчення технічних засобів їхнього усунення. Обов'язкове вивчення цих правил спрямоване на мінімізацію ризиків, пов'язаних із соціальною інженерією та фішингом.

Крім того слід зауважити, що при вивченні змістової лінії ``Мультимедійні та гіпертекстові документи'', тема створення веб-сайтів за допомогою систем керування вмістом включає вимогу дотримання правил ергономічного розміщення матеріалів. Хоча це безпосередньо стосується використання та зручності для користувача, всеж таки це є елементом забезпечення надійності та цілісності ресурсу, що опосередковано стосується безпеки, оскільки добре структурований і ергономічний сайт менш схильний до помилок, які можуть створити вразливості, але, знову ж таки, технічний захист сервера чи коду від злому тут не розглядається.

Обов'язкова частина, представлена базовим модулем, охоплює лише загальні основи інформаційної безпеки та її проблеми у широкому соціальному контексті. Вона повністю фокусується на етичній поведінці та цифровому громадянстві як головних інструментах захисту, що включає усвідомлення загроз в Інтернеті. Навички базового модуля охоплюють загальну комп'ютерну грамотність та планування діяльності, тоді як технічна безпека залишається виключно на рівні правил поведінки. У цій частині повністю відсутні теми мережевого захисту та збереження даних на системному рівні, оскільки вони виходять за рамки мінімальних вимог.

Базовий модуль ефективно інтегрує деякі основи безпеки, закладаючи необхідний фундамент для вибіркового модулю ``Інформаційна безпека''.


\subsubsection{Аналіз вибіркового модуля ``Інформаційна безпека''} 
Програма вибіркового модулю ``Інформаційна безпека'' рекомендована Міністерством освіти і науки України \cite{inform_10_11_standart}. Вона рекомендована для учнів профільної школи і розрахована на 17 годин. Програма охоплює три розділи.

Розділ 1 ``Основи безпеки інформаційних технологій'' присвячений формуванню ціннісної основи кібербезпеки, а також правової та аналітичної свідомості. Під час вивчення розділу учні опановують основні поняття інформаційної безпеки, зрозуміють місце та роль автоматизованих систем, а також причини загострення проблеми безпеки. Детально вивчають інформаційні відносини, суб'єктів, їхні інтереси та шляхи нанесення шкоди. Аналізують загрози безпеці, джерела та канали несанкціонованого доступу, включаючи комп'ютерні віруси та програми-зломщики. Основний акцент зміщується з базових навичок (які учні вже мають після базового модулю курсу інформатики) на правове обґрунтування, систематизацію загроз та аналіз ризиків. Учні осмислюють важливість і місце інформаційної безпеки в сучасному суспільстві. Учні повинні навчитися не лише дотримуватись правил кібергігієни, але й аналізувати вразливості, створювати базові матриці ризиків для інформаційних систем, а також розуміти правове поле інформаційних відносин та відповідальності в Україні. 
Наприкінці вивчення цього розділу учні зможуть не лише ідентифікувати базові загрози, але й сформує правове та етичне підґрунтя для подальшої практичної роботи з безпекою, усвідомлюючи власну відповідальність у цифровому середовищі.

Розділ 2 ``Забезпечення безпеки інформаційних технологій'' є фундаментальним і присвячений системному захисту окремого комп'ютера та даних. Його зміст є найбільш насиченим з погляду нових, складних тем, які не розглядалися на базовому рівні, зокрема криптографія та моделі розмежування доступу. Це закладає основи інженерного підходу до захисту інформації. Під час вивчення розділу учні детально опанують керування доступом: моделі ІАА (ідентифікація, автентифікація, авторизація), дискреційну та мандатну моделі. Проведуть детальне вивчення криптографічних методів (шифрування та хешування), їх класифікації. Ознайомляться з електронним цифровим підписом (ЕЦП) та його інфраструктурою (центри сертифікації). Розглянуть засоби контролю цілісності програмного забезпечення та даних, а також навчаться створювати і відновлювати резервні копії операційних систем і даних, забезпечуючи цілісність та конфіденційність інформації. 
Результати цього розділу забезпечують глибоке розуміння та практичне володіння фундаментальними інструментами захисту даних і систем, зокрема криптографічними механізмами, що є наріжним каменем будь-якої інформаційної безпеки.

Розділ 3 ``Забезпечення безпеки комп'ютерних систем і мереж'' найбільш практико-орієнтований. Він повністю присвячений мережевому захисту та конфігуруванню інфраструктури. Опанування матеріалом роздулу сприятиме підготовці учнів до роботи з корпоративними мережами. Учні навчаться аналізувати типову корпоративну мережу, рівні інфраструктури, мережеві загрози (DDoS, MITM) та атаки. Опанують віртуальні приватні мережі (VPN), їх призначення та тунелювання. Вивчать призначення та захисні механізми міжмережевих екранів (брандмауерів), а також політику безпеки.  Розглянуть системи аналізу вмісту трафіку (пошта, веб) та сценарії реагування.
Теми міжмережеві екрани (брандмауери), VPN-мережі та аналіз трафіку є абсолютно новими для учнів. При вивченні розділу учні повинні набути діяльнісної складової у сфері захисту мереж. Вони мають вміти аналізувати мережеві загрози (DDoS, MITM), розуміти принципи роботи брандмауерів та конфігурувати правила фільтрації (на рівні симуляції чи ОС). Крім того, вони повинні вміти налаштовувати VPN-клієнтів та розуміти, як працюють системи аналізу вмісту трафіку. Результатом є практичні навички захисту мережевої інфраструктури. 
Після завершення розділу учні мають пояснити типи корпоративних мереж, класифікувати загрози, описувати принципи роботи брандмауерів та VPN. Виконувати конфігурування міжмережевого екрану, розуміючи принципи функціонування корпоративних мереж та вміючи конфігурувати критичні елементи захисту, такі як брандмауери та VPN. Також використовувати засоби моніторингу мережного трафіку, виконувати конфігурування простих маршрутизаторів. Набути  практичні навички роботи з периметровим захистом, усвідомлювати необхідність та виконувати резервне збереження даних, дотримуватись правил безпечної роботи в Інтернеті, враховувати наслідки несанкціонованого доступу,
 

 Аналіз програми дозволив запропонувати такий календарний план (таблиця~\ref{tab:plan}).

{\small
\setlength \tabcolsep {0.4\tabcolsep }
\begin{longtable} {|c|p{0.3\linewidth}|p{0.3\linewidth}|p{0.3\linewidth}|}
\caption{\normalsize Календарно-тематичне планування для вибіркового модуля ``Інформаційна безпека''}

\label{tab:plan}\\
\hline
\textbf{№} & \hfil \textbf{Тема уроку} & \hfil \textbf{Мета уроку} & \hfil \textbf{Зміст навчання} \\
\hline
\endfirsthead % Заголовок для першої сторінки

\multicolumn{4}{r}{{\normalsize Продовження табл. \thetable}} \\
\hline
\endhead % Заголовок для наступних сторінок


\multicolumn{4}{|c|}{\textbf{Розділ I: Основи безпеки ІТ (систематизація) (3 год.)}} \\
\hline
1 & Систематизація загроз. Правове поле ІБ в Україні. & Систематизувати знання про загрози та вразливості; засвоїти правові основи ІБ (кейс-стаді). & Теоретичні основи. Правове поле ІБ в Україні. \\
\hline
2 & Інформаційні відносини та аналіз вразливостей. Оцінка ризиків. Практична робота №1 ``Систематизація загроз та правове поле інформаційної безпеки'' & Опанувати методологію оцінки ризиків; практично створити матрицю ризиків. & Інформаційні відносини. Аналіз вразливостей. Практикум: Створення матриці ризиків. \\
\hline
3 & Політика кібергігієни (переосмислення) та контроль цілісності (вступ). & Навчитися використовувати хеш-суми для контролю цілісності; переосмислити правила кібергігієни. & Політика кібергігієни. Практикум: використання хеш-сум. \\
\hline
\multicolumn{4}{|c|}{\textbf{Розділ II: Системний захист (криптографія)  (7 год.)}} \\
\hline
4 & Керування доступом: ІАА. Дискреційна/мандатна моделі. & Засвоїти ключові моделі керування доступом (ІАА); відпрацювати застосування через рольову гру. & Керування доступом: ідентифікація, автентифікація, авторизація (ІАА). Моделі доступу. \\
\hline
5 & Криптографічні методи: Шифрування та хешування (теорія). & Зрозуміти принципи та класифікацію криптографічних методів; проаналізувати алгоритми. & Криптографічні методи: шифрування (симетричне/асиметричне) та хешування. \\
\hline
6 & Практична робота №2 ``Застосування шифрування на практицi'' & Набути практичних навичок шифрування та дешифрування із застосуванням GPG / аналогів. & Діяльнісна складова: робота з GPG / аналогами (симетричне та асиметричне шифрування). \\
\hline
7 & Електронний цифровий підпис (ЕЦП). Центри сертифікації. & Вивчити принципи ЕЦП та інфраструктуру відкритих ключів; проаналізувати приклади використання. & ЕЦП та інфраструктура відкритих ключів. Кейс: Використання ЕЦП в Україні. \\
\hline
8 & Засоби контролю цілісності ПЗ та даних. & Ознайомитися з інструментами контролю цілісності; практично перевірити цілісність ПЗ. & Засоби контролю цілісності ПЗ та даних. Практикум: Перевірка цілісності ПЗ. \\
\hline
9 & Резервне копіювання ОС та даних. Політика відновлення. & Розробити політику резервного копіювання та відновлення; практично створити образ ОС (віртуальне середовище). & Резервне копіювання ОС та даних. Практикум: Створення образу ОС. \\
\hline
\multicolumn{4}{|c|}{\textbf{Розділ III: Мережевий захист (конфігурування) (7 год.)}} \\
\hline
11 & Корпоративна мережа. Рівні інфраструктури. Мережеві загрози. & Вивчити структуру корпоративної мережі та проаналізувати основні мережеві загрози (DDoS, MITM). & Корпоративна мережа, її інфраструктура. Мережеві загрози (DDoS, MITM). \\
\hline
12 & Міжмережеві екрани (брандмауери): принципи роботи. & Засвоїти принципи функціонування брандмауерів; зрозуміти їхню роль у периметровому захисті. & Міжмережеві екрани (брандмауери): принципи роботи та захисні механізми. \\
\hline
13 & Практична робота 3 ``Конфігурування міжмережевого екрану''  & Набути практичних навичок конфігурування брандмауера та налаштування правил фільтрації. & Діяльнісна складова: налаштування правил фільтрації (ОС/симулятор). \\
\hline
14 & Віртуальні приватні мережі (VPN): призначення та тунелювання. & Вивчити призначення та архітектуру VPN; зрозуміти механізми тунелювання. & Віртуальні приватні мережі (VPN): призначення та тунелювання. \\
\hline
15 & Практична робота 4 ``Налаштування VPN-клієнта''. & Практично налаштувати та використовувати VPN-клієнт для безпечного віддаленого доступу. & Діяльнісна складова: робота з VPN-клієнтом. \\
\hline
16 & Системи аналізу вмісту трафіку (пошта, веб). & Ознайомитися з принципами роботи систем DLP/IDS/IPS та їхньою роллю у фільтрації контенту. & Системи аналізу вмісту трафіку (пошта, веб). \\
\hline
17 & Підсумкове заняття / Захист проєкту. & Систематизувати набуті знання та вміння; здійснити презентацію та захист фінального проєкту. & Систематизація знань, оцінювання, захист проєкту. \\
\hline
\end{longtable}
}


\subsection{Розробка завдань для диференційованого навчання вибіркового модуля ``Інформаційна безпека'' учнів ліцею}





\subsubsection{Практична робота №1 ``Систематизація загроз та правове поле інформаційної безпеки''}

Розроблена практична робота демонструє реалізацію диференційованого підходу при вивченні першої теми курсу. Структурно робота організована таким чином, що забезпечує гнучкість у виборі траєкторії виконання завдань при збереженні єдиних навчальних цілей.

Практична робота виконується протягом двох уроків з додатковою самостійною роботою вдома. Перший урок (45 хвилин) присвячений виконанню завдання середнього рівня для кожного спрямування в класі під керівництвом учителя. Це забезпечує формування фундаментальних знань та вмінь, необхідних для подальшої самостійної роботи. Завдання достатнього та високого рівнів виконується як домашня робота за вибором учня, що дозволяє індивідуалізувати темп та глибину опрацювання матеріалу. Презентація результатів (15 хвилин на наступному уроці) створює умови для обміну досвідом та взаємонавчання учнів.

Академічний напрямок ``Дослідник безпеки'' спрямований на формування аналітичних та дослідницьких компетентностей у галузі інформаційної безпеки. Завдання цього напрямку (табл. \ref{tab:lab1_acad_task0}) акцентують увагу на систематизації знань, критичному аналізі джерел та створенні інформаційних продуктів .

% Визначимо колір для заголовка таблиці
\definecolor{TableHeaderBlue}{RGB}{46, 82, 141}

% -- Початок таблиці ---
% \noindent запобігає відступу, а [!] просить LaTeX ``спробувати''
% розмістити таблицю тут.
\begin{table}[h!]
\caption{Завдання для ``Дослідників безпеки''.}
\label{tab:lab1_acad_task0} % Мітка для посилань на таблицю
\noindent % Забирає відступ зліва
% Налаштування для списків всередині таблиці:
% nosep - прибирає зайві вертикальні відступи
% leftmargin=* - стандартний відступ для списку
\setlist[enumerate]{nosep, leftmargin=*}


% p{width} - колонка з фіксованою шириною та авто-переносом
% >{\raggedright\arraybackslash} - вирівнювання по лівому краю (краще для списків)
\begin{tabular}{
  | >{\raggedright\arraybackslash}p{0.15\textwidth} 
  | >{\raggedright\arraybackslash}p{0.80\textwidth} |
}
\hline

% -- Кольоровий заголовок ---
% \rowcolor - колір фону, \color - колір тексту
\rowcolor{TableHeaderBlue}
\hfil \color{white}\textbf{РІВЕНЬ} & 
\hfil \color{white}\textbf{ЗАВДАННЯ} \\ \hline

% -- Рівень I ---
\textbf{Середній} (6 балів)&
\textbf{Базове дослідження:}
\begin{enumerate}
    \item Створити таблицю-класифікацію з 8+ загроз ІБ за категоріями: джерело (внутрішні/зовнішні), намір (навмисні/випадкові), об’єкт атаки
    \item Знайти 2 реальні приклади кіберінцидентів в Україні (2023-2024 рр.)
    \item Оформити у вигляді документа з посиланнями на джерела
\end{enumerate} \\ \hline

% -- Рівень II ---
\textbf{Достатній} (9 балів)&
\textbf{Поглиблений аналіз (додатково до середнього рівня):}
\begin{enumerate}
    \item Побудувати схему взаємозв’язків між типами загроз
    \item Порівняти класифікації загроз: Закон України ``Про основні засади забезпечення кібербезпеки'' vs ISO 27001
    \item Проаналізувати тенденції змін загроз за останні 3 роки
\end{enumerate} \\ \hline

% -- Рівень III ---
\textbf{Високий} (12 балів)&
\textbf{Експертний рівень (додатково до достатнього рівня):}
\begin{enumerate}
    \item Створити інфографіку ``Ландшафт загроз ІБ для освітніх закладів України''
    \item Підготувати презентацію до виступу на тему ``Як змінилися загрози ІБ в Україні після 24.02.2022''
    \item Розробити 3 рекомендації для ліцею щодо протидії актуальним загрозам
\end{enumerate} \\ \hline
\end{tabular}
\end{table}
% -- Кінець таблиці ---


На середньому рівні учні створюють структуровану класифікацію загроз інформаційній безпеці, використовуючи критерії джерела походження (внутрішні/зовнішні), наміру (навмисні/випадкові) та об'єкта атаки. Обов'язковим компонентом є пошук та аналіз двох реальних прикладів кіберінцидентів в Україні за останні роки. Це формує навички роботи з актуальною інформацією та розуміння реальних загроз у національному контексті.

Достатній рівень передбачає побудову схеми взаємозв'язків між типами загроз, що розвиває системне мислення учнів. Порівняльний аналіз класифікацій загроз за українським законодавством та міжнародним стандартом ISO 27001 формує розуміння різних підходів до систематизації загроз. Аналіз тенденцій змін загроз за останні три роки розвиває навички прогнозування та стратегічного мислення.

На високому рівні учні створюють комплексні інформаційні продукти~-- інфографіку ландшафту загроз для освітніх закладів України та презентацію про трансформацію загроз після початку повномасштабної війни. Розробка практичних рекомендацій для закладу освіти демонструє здатність застосовувати теоретичні знання для вирішення реальних проблем.


Технічний напрямок ``Етичний хакер'' орієнтований на формування практичних навичок виявлення та усунення вразливостей інформаційних систем. Завдання (табл. \ref{tab:lab1_tech_task}) побудовані за принципом поступового ускладнення від базової діагностики до симуляції пентесту.

\definecolor{TableHeaderBlue}{RGB}{46, 82, 141}
\definecolor{TableBodyGreen}{RGB}{229, 248, 230}
\begin{table}[h!]
\caption{Завдання для ``Етичного хакера''.}
\label{tab:lab1_tech_task} % Мітка для посилань на таблицю
\noindent % Забирає відступ зліва
% Налаштування для списків всередині таблиці:
% nosep - прибирає зайві вертикальні відступи
% leftmargin=* - стандартний відступ для списку
\setlist[enumerate]{nosep, leftmargin=*}

% p{width} - колонка з фіксованою шириною та авто-переносом
% >{\raggedright\arraybackslash} - вирівнювання по лівому краю (краще для списків)
\begin{tabular}{
  | >{\raggedright\arraybackslash}p{0.15\textwidth} 
  | >{\raggedright\arraybackslash}p{0.8\textwidth} |
}
\hline

% -- Кольоровий заголовок ---
% \rowcolor - колір фону, \color - колір тексту
\rowcolor{TableHeaderBlue}
\hfil\color{white}\textbf{РІВЕНЬ} & 
\hfil\color{white}\textbf{ЗАВДАННЯ} \\ \hline

% -- Рівень I ---
% \cellcolor - встановлює колір для окремої комірки
\cellcolor{TableBodyGreen}
\textbf{Середній} (6 балів) &
\textbf{Базова практика:}
\begin{enumerate}
    \item Провести сканування безпеки ПК (Shields Up! від GRC.com)
    \item Заповнити чек-лист: антивірус, оновлення ОС, складність паролів, фаєрвол
    \item Створити скріншот-звіт з поясненнями знайдених проблем
\end{enumerate} \\ \hline

% -- Рівень II ---
\cellcolor{TableBodyGreen}
\textbf{Достатній} (9 балів) &
\textbf{Технічне дослідження (додатково до рівня І):}
\begin{enumerate}
    \item Дослідити одну CVE вразливість: опис, механізм, патч
    \item Створити демонстрацію загрози: фішинг (макет), слабкий пароль, або SQL-ін’єкція (DVWA)
    \item Протестувати з онлайн-сервіси на HaveIBeenPwned
\end{enumerate} \\ \hline

% -- Рівень III ---
\cellcolor{TableBodyGreen}
\textbf{Високий} (12 балів) &
\textbf{Пентест-симуляція (додатково до рівнів І-ІІ):}
\begin{enumerate}
    \item Провести міні-аудит за OWASP Top 10
    \item Написати Python-скрипт для перевірки складності паролів
    \item Записати відео-демонстрацію (3-5 хв) атаки та захисту
\end{enumerate} \\ \hline

\end{tabular}
\end{table}
% -- Кінець таблиці ---


Базовий рівень включає проведення сканування безпеки персонального комп'ютера з використанням онлайн-інструментів, зокрема Shields Up! від GRC.com. Учні опановують базові навички технічного аудиту через заповнення чек-листа, що включає перевірку антивірусного захисту, актуальності оновлень операційної системи, складності паролів та налаштувань фаєрволу. Документування результатів у вигляді скріншот-звіту з поясненнями формує навички технічної документації.

На технічному рівні дослідження учні поглиблюють розуміння механізмів вразливостей через дослідження конкретної CVE-вразливості, включаючи її опис, механізм експлуатації та методи усунення. Створення демонстрації загрози (макет фішингової атаки, демонстрація атаки на слабкий пароль або SQL-ін'єкція в навчальному середовищі DVWA) забезпечує практичне розуміння методів атак. Використання сервісу HaveIBeenPwned для перевірки витоків даних формує навички персональної кібергігієни.

Рівень пентест-симуляції передбачає проведення міні-аудиту за методологією OWASP Top 10, що є професійним стандартом у галузі веб-безпеки. Написання Python-скрипту для перевірки складності паролів розвиває навички програмування для вирішення задач безпеки. Створення відео-демонстрації атаки та захисту (3-5 хвилин) формує вміння візуалізувати та пояснювати складні технічні процеси.

Управлінський напрямок ``Менеджер безпеки'' спрямований на формування компетентностей з управління інформаційною безпекою на організаційному рівні. Завдання (табл. \ref{tab:lab1_manag_task}) акцентують увагу на правових, економічних та організаційних аспектах забезпечення інформаційної безпеки.


\definecolor{TableHeaderBrown}{RGB}{142, 85, 34}
\definecolor{TableBodyPeach}{RGB}{254, 237, 203}
\begin{table}[h!]
\caption{Завдання для ``Менеджера безпеки''.}
\label{tab:lab1_manag_task} % Мітка для посилань на таблицю
\noindent % Забирає відступ зліва
% Налаштування для списків всередині таблиці
\setlist[enumerate]{nosep, leftmargin=*}

\begin{tabular}{
  | >{\raggedright\arraybackslash}p{0.15\textwidth} 
  | >{\raggedright\arraybackslash}p{0.8\textwidth} |
}
\hline

% -- Кольоровий заголовок ---
\rowcolor{TableHeaderBrown}
\hfil \color{white}\textbf{РІВЕНЬ} & 
\hfil \color{white}\textbf{ЗАВДАННЯ} \\ \hline

% -- Рівень I ---
\cellcolor{TableBodyPeach}
\textbf{Середній} (6 балів) &
\textbf{Базове планування:}
\begin{enumerate}
    \item Проаналізувати Закон України ``Про захист персональних даних''
    \item Розробити чек-лист compliance для інтернет-магазину
    \item Створити шаблон ``Згоди на обробку персональних даних''
\end{enumerate} \\ \hline

% -- Рівень II ---
\cellcolor{TableBodyPeach}
\textbf{Достатній} (9 балів) &
\textbf{Управлінське рішення (додатково до Рівня І):}
\begin{enumerate}
    \item Розробити матрицю ризиків 5х5 для стартапу (10 загроз)
    \item Розрахувати базовий бюджет ІБ для компанії з 10 співробітників
    \item Створити план реагування на інцидент (1 сторінка)
\end{enumerate} \\ \hline

% -- Рівень III ---
\cellcolor{TableBodyPeach}
\textbf{Високий} (12 балів) &
\textbf{Стратегічний рівень (додатково до Рівнів І-ІІ):}
\begin{enumerate}
    \item Розробити презентацію для керівництва (7-10 слайдів) з ROI
    \item Створити проект ``Положення про інформаційну безпеку'' (2-3 сторінки)
    \item Розробити 5 KPI для оцінки ефективності ІБ
\end{enumerate} \\ \hline

\end{tabular}
\end{table}
% -- Кінець таблиці ---

Базовий рівень планування включає аналіз Закону України ``Про захист персональних даних'' та розробку практичних документів для забезпечення відповідності вимогам законодавства. Створення чек-листа compliance для інтернет-магазину та шаблону згоди на обробку персональних даних формує розуміння правових вимог та навички розробки корпоративної документації.

Управлінське рішення передбачає розробку матриці ризиків розміром 5×5 для стартапу з ідентифікацією та оцінкою 10 загроз. Розрахунок базового бюджету інформаційної безпеки для малого підприємства розвиває навички фінансового планування в галузі безпеки. Створення плану реагування на інцидент формує розуміння процесів управління інцидентами.

На стратегічному рівні учні розробляють презентацію для керівництва з обґрунтуванням повернення інвестицій (ROI) в інформаційну безпеку. Створення проекту ``Положення про інформаційну безпеку'' демонструє здатність розробляти нормативні документи організації. Розробка ключових показників ефективності (KPI) для оцінки стану інформаційної безпеки формує навички стратегічного управління.

\subsubsection{Урок №4 ``Керування доступом: ІАА. Дискреційна/Мандатна моделі''}

Тема керування доступом є фундаментальною у вивченні інформаційної безпеки, оскільки формує розуміння базових механізмів захисту інформаційних ресурсів. Розроблена система диференційованих завдань для уроку №4 реалізує комплексний підхід до вивчення концепції ІАА (Ідентифікація, Автентифікація, Авторизація) та моделей керування доступом через призму професійних інтересів та майбутньої спеціалізації учнів.

Структура уроку побудована за принципом поступового занурення у предметну область: від загального розуміння процесів ідентифікації користувачів до практичної реалізації різних моделей керування доступом. Особливістю методичного підходу є використання спільної стартової активності, що забезпечує формування єдиного понятійного апарату перед диференційованою роботою.

Мозковий штурм ``Як система впізнає користувача?'' створює когнітивну основу для подальшої диференційованої роботи. Через перегляд демонстраційного відео процесу входу в систему (2 хвилини) та інтерактивне заповнення схеми (рис.~\ref{fig:koncepcia_iaa}) на дошці учні формують розуміння триєдиної концепції ІАА. Схема візуалізує логічний ланцюжок: ``Хто я?'' відповідає процесу ідентифікації через логін, ``Доведи це!'' розкриває сутність автентифікації через паролі чи біометрію, ``Що можу?'' демонструє механізм авторизації через права доступу.

\begin{figure}[h!]
\centering
\begin{tikzpicture}[
    node distance=18mm,
    box/.style={
        rectangle,
        draw,
        rounded corners,
        text width=12cm,
        align=left,
        fill=gray!10,
        drop shadow
    },
    arrow/.style={-{Latex[length=3mm]}, thick}
]

% Title
\node[box, text width=12cm, align=center] (title) { \textbf{ПРОЦЕС ВХОДУ В СИСТЕМУ} };

% Identification
\node[box, below=of title] (id) {
\textbf{1. ІДЕНТИФІКАЦІЯ — ``Хто я?''}\\
\hspace*{1em}• Введення логіну (username)\\
\hspace*{1em}• Система знаходить обліковий запис
};

% Authentication
\node[box, below=of id] (auth) {
\textbf{2. АВТЕНТИФІКАЦІЯ — ``Доведи це!''}\\
\hspace*{1em}• Пароль (щось, що знаєш)\\
\hspace*{1em}• Токен/картка (щось, що маєш)\\
\hspace*{1em}• Біометрія (щось, чим є)
};

% Authorization
\node[box, below=of auth] (author) {
\textbf{3. АВТОРИЗАЦІЯ — ``Що можу робити?''}\\
\hspace*{1em}• Перевірка прав доступу\\
\hspace*{1em}• Застосування політик безпеки\\
\hspace*{1em}• Надання дозволів на ресурси
};

% Arrows
\draw[arrow] (title) -- (id);
\draw[arrow] (id) -- (auth);
\draw[arrow] (auth) -- (author);

\end{tikzpicture}
\caption{Концепція IAA (ідентифікація, автентифікація, авторизація).}
\label{fig:koncepcia_iaa}
\end{figure}


Така структурована візуалізація забезпечує формування ментальної моделі, яка стає опорою для виконання диференційованих завдань незалежно від обраного напрямку. Важливо, що всі учні отримують однакову базову інформацію, але далі інтерпретують та застосовують її відповідно до специфіки свого профілю.



% Визначимо колір для заголовка таблиці
\definecolor{TableHeaderBlue}{RGB}{46, 82, 141}

% -- Початок таблиці ---
% \noindent запобігає відступу, а [!] просить LaTeX ``спробувати''
% розмістити таблицю тут.
\begin{table}[h!]
\caption{Завдання для академічного напрямку ``Аналітик систем безпеки''.}
\label{tab:lesson4_acad_task} % Мітка для посилань на таблицю
\noindent % Забирає відступ зліва
% Налаштування для списків всередині таблиці:
% nosep - прибирає зайві вертикальні відступи
% leftmargin=* - стандартний відступ для списку
\setlist[enumerate]{nosep, leftmargin=*}


% p{width} - колонка з фіксованою шириною та авто-переносом
% >{\raggedright\arraybackslash} - вирівнювання по лівому краю (краще для списків)
\begin{tabular}{
  | >{\raggedright\arraybackslash}p{0.15\textwidth} 
  | >{\raggedright\arraybackslash}p{0.80\textwidth} |
}
\hline

% -- Кольоровий заголовок ---
% \rowcolor - колір фону, \color - колір тексту
\rowcolor{TableHeaderBlue}
\hfil \color{white}\textbf{РІВЕНЬ} & 
\hfil \color{white}\textbf{ЗАВДАННЯ} \\ \hline

\textbf{Середній} (6 балів) &
\textbf{Базове дослідження:}
\begin{enumerate}
    \item Заповніть порівняльну таблицю Дискреційної (DAC) і Мандатної (MAC) моделей керування доступом за критеріями: принцип роботи, переваги, недоліки, приклади застосування
    \item Знайдіть 2 реальні приклади використання IAA в українських онлайн-сервісах (Дія, банкінг тощо)
    \item Оформити у вигляді аналітичної записки з посиланнями на джерела
\end{enumerate} \\ \hline

\textbf{Достатній} (9 балів)&
\textbf{Поглиблений аналіз (додатково до середнього рівня):}
\begin{enumerate}
    \item Порівняти реалізацію автентифікації: однофакторна vs багатофакторна (MFA)
    \item Проаналізувати модель RBAC (Role-Based Access Control) та її застосування в освітніх закладах
    \item Створити схему взаємозв'язків між компонентами IAA
\end{enumerate} \\ \hline

\textbf{Високий} (12 балів)&
\textbf{Експертний рівень (додатково до достатнього рівня):}
\begin{enumerate}
    \item Підготувати порівняльний аналіз моделей керування доступом (DAC, MAC, RBAC, ABAC) з рекомендаціями щодо застосування
    \item Розробити проєкт системи керування доступом для шкільної мережі
    \item Створити презентацію ``Сучасні методи автентифікації: від паролів до біометрії''
\end{enumerate} \\ \hline
\end{tabular}
\end{table}
% -- Кінець таблиці ---

Академічний напрямок ``Аналітик систем безпеки'' (табл.~\ref{tab:lesson4_acad_task}) орієнтований на формування аналітичних компетентностей у сфері керування доступом. На середньому рівні учні опановують базові концепції IAA та порівнюють основні моделі керування доступом на прикладах реальних систем. Достатній рівень передбачає поглиблений аналіз механізмів автентифікації та авторизації, включаючи сучасні підходи RBAC. На високому рівні учні демонструють здатність синтезувати знання для проєктування систем керування доступом та презентації результатів дослідження.

Технічний напрямок ``Системний адміністратор'' (табл.~\ref{tab:lesson4_tech_task}) фокусується на практичних навичках налаштування прав доступу в операційній системі Linux, що є стандартом де-факто для серверних систем. Завдання побудовані за принципом «від команди до автоматизації».



\definecolor{TableHeaderBlue}{RGB}{46, 82, 141}
\definecolor{TableBodyGreen}{RGB}{229, 248, 230}
\begin{table}[h!]
\caption{Завдання для ``Системного адміністратора''.}
\label{tab:lesson4_tech_task} % Мітка для посилань на таблицю
\noindent % Забирає відступ зліва
% Налаштування для списків всередині таблиці:
% nosep - прибирає зайві вертикальні відступи
% leftmargin=* - стандартний відступ для списку
\setlist[enumerate]{nosep, leftmargin=*}

% p{width} - колонка з фіксованою шириною та авто-переносом
% >{\raggedright\arraybackslash} - вирівнювання по лівому краю (краще для списків)
\begin{tabular}{
  | >{\raggedright\arraybackslash}p{0.15\textwidth} 
  | >{\raggedright\arraybackslash}p{0.8\textwidth} |
}
\hline

% -- Кольоровий заголовок ---
% \rowcolor - колір фону, \color - колір тексту
\rowcolor{TableHeaderBlue}
\hfil\color{white}\textbf{РІВЕНЬ} & 
\hfil\color{white}\textbf{ЗАВДАННЯ} \\ \hline

% -- Рівень I ---
% \cellcolor - встановлює колір для окремої комірки
\cellcolor{TableBodyGreen}
\textbf{Середній} (6 балів) &
\textbf{Базова практика:}
\begin{enumerate}
    \item Виконати команди Linux:
    \begin{itemize}
        \item Створення користувачів: \texttt{alice}, \texttt{bob}, \texttt{charlie}. \vspace{-0.5em}
        \item Створення файлів та встановлення прав: \texttt{chmod 600}, \texttt{chmod 644}. \vspace{-0.5em}
    \end{itemize}
    \item Пояснити значення чисел у \texttt{chmod 644}.
    \item Міні-тест: намалювати права для \texttt{chmod 755} (rwxr-xr-x).
\end{enumerate}
\\ \hline

% -- Рівень II ---
\cellcolor{TableBodyGreen}
\textbf{Достатній} (9 балів) &
\textbf{Налаштування груп:}
\begin{enumerate}
    \item Створити групи: \texttt{developers}, \texttt{managers}; додати користувачів.
    \item Реалізувати доступ:
    \begin{itemize}
        \item /project — доступ лише developers. \vspace{-0.5em}
        \item /reports — managers читають, лише bob пише. \vspace{-0.5em}
        \item /public — всі читають, ніхто не пише. \vspace{-0.5em}
    \end{itemize}
    \item Знайти помилку: файл \texttt{password.txt} має права \texttt{rwxrwxrwx}. Чому це небезпечно? Як виправити?
\end{enumerate} \\ \hline

% -- Рівень III ---
\cellcolor{TableBodyGreen}
\textbf{Високий} (12 балів) &
\textbf{Скриптування :}
\begin{enumerate}
    \item Написати bash-скрипт, що реалізує:
    \begin{itemize}
        \item створення користувача, \vspace{-0.5em}
        \item додавання до групи, \vspace{-0.5em}
        \item створення домашньої директорії з правами 700, \vspace{-0.5em}
        \item логування дій. \vspace{-0.5em}
    \end{itemize}
    \item Налаштувати правила sudoers:
    \begin{itemize}
        \item developers можуть перезапускати вебсервер, \vspace{-0.5em}
        \item managers можуть переглядати логи, \vspace{-0.5em}
        \item ніхто не може видаляти системні файли. \vspace{-0.5em}
    \end{itemize}
    \item Домашнє завдання: встановити VirtualBox, створити Linux VM, реалізувати структуру.
\end{enumerate} \\ \hline

\end{tabular}
\end{table}
% -- Кінець таблиці ---



На середньому рівні учні виконують базові операції: створення користувачів (useradd alice, bob, charlie), створення файлів з різними правами доступу, розуміння октальної нотації прав (chmod 600, 644).
Декодування числової нотації прав формує глибоке розуміння механізму: 6 означає читання+запис для власника (4+2), 4 - тільки читання для групи та інших. Міні-тест з інтерпретації chmod 755 та візуалізації у форматі rwxr-xr-x закріплює розуміння відповідності між числовою та символьною нотацією.

Достатній рівень вводить концепцію груп як механізму організації доступу. Створення структури груп (developers, managers) та призначення користувачів демонструє практичну реалізацію рольової моделі в Linux. Реалізація складного сценарію з трьома папками (/project для developers, /reports для managers з диференційованими правами, /public для загального читання) формує навички проектування файлової структури.
Критично важливим є завдання з пошуку та виправлення помилки безпеки - файл password.txt з правами rwxrwxrwx (777) під власністю root. Це формує розуміння критичності правильного налаштування прав для чутливих даних та навички аудиту безпеки.

Високий рівень передбачає написання bash-скрипту для автоматизації процесу створення користувачів з правильними правами. Скрипт повинен включати інтерактивне введення даних, логіку визначення департаменту, автоматичне додавання до груп, створення домашньої директорії з обмеженими правами (700) та логування всіх дій для аудиту.
Реалізація sudo-правил через файл sudoers демонструє механізм делегування привілеїв: developers отримують право перезапуску веб-сервера без повних прав root, managers можуть переглядати логи для моніторингу, але ніхто не може видаляти системні файли. Це формує розуміння принципу найменших привілеїв на практичному рівні.
Домашнє завдання з встановлення VirtualBox та створення Linux VM переводить навчання на новий рівень - від симуляторів до реального середовища, підготовлюючи учнів до професійної діяльності.


%Управлінський напрямок ``Менеджер безпеки'' спрямований на формування компетентностей з управління інформаційною безпекою на організаційному рівні. Завдання (табл. \ref{tab:lab1_manag_task}) акцентують увагу на правових, економічних та організаційних аспектах забезпечення інформаційної безпеки.

Управлінський напрямок орієнтований на формування компетентностей розробки організаційних політик та процедур керування доступом. Завдання (табл. \ref{tab:lesson4_manag_task}) моделюють реальні управлінські ситуації від малого офісу до корпоративного рівня.

\definecolor{TableHeaderBrown}{RGB}{142, 85, 34}
\definecolor{TableBodyPeach}{RGB}{254, 237, 203}
\begin{table}[h!]
\caption{Завдання для ``Менеджера безпеки''.}
\label{tab:lesson4_manag_task} % Мітка для посилань на таблицю
\noindent % Забирає відступ зліва
% Налаштування для списків всередині таблиці
\setlist[enumerate]{nosep, leftmargin=*}

\begin{tabular}{
  | >{\raggedright\arraybackslash}p{0.15\textwidth} 
  | >{\raggedright\arraybackslash}p{0.8\textwidth} |
}
\hline

% -- Кольоровий заголовок ---
\rowcolor{TableHeaderBrown}
\hfil \color{white}\textbf{РІВЕНЬ} & 
\hfil \color{white}\textbf{ЗАВДАННЯ} \\ \hline

% -- Рівень I ---
\cellcolor{TableBodyPeach}
\textbf{Середній} (6 балів) &
\textbf{Розробка базової політики}
\begin{enumerate}
    \item Заповнити шаблон політики доступу (малий офіс, 10 осіб):
    \begin{itemize}
        \item рівні доступу: публічний / внутрішній / конфіденційний, \vspace{-0.5em}
        \item ролі: директор, бухгалтер, менеджер, працівник,\vspace{-0.5em}
        \item правила паролів: довжина, період зміни, кількість спроб.
    \end{itemize}
    \item Міні-кейс: кроки надання доступу новому працівнику.
\end{enumerate} \\ \hline

% -- Рівень II ---
\cellcolor{TableBodyPeach}
\textbf{Достатній} (9 балів) &
\textbf{Комплексна система управління :}
\begin{enumerate}
    \item Розробити матрицю RACI (директор, ІТ, HR, керівник відділу).
    \item Створити чек-лист деактивації доступу (8--10 пунктів).
    \item Оцінити ризики: три сценарії, якщо доступ не забрати вчасно.
\end{enumerate} \\ \hline

% -- Рівень III ---
\cellcolor{TableBodyPeach}
\textbf{Високий} (12 балів) &
\textbf{Стратегічне планування системи:}
\begin{enumerate}
    \item Розробити систему Identity Management для компанії 100+ осіб:
    \begin{itemize}
        \item Active Directory,\vspace{-0.5em}
        \item SSO,\vspace{-0.5em}
        \item onboarding/offboarding,\vspace{-0.5em}
        \item KPI.
    \end{itemize}
    \item Скласти бюджет впровадження: ліцензії, навчання, аудит, підтримка, ROI.
    \item Створити презентацію (5 слайдів): проблеми, рішення, переваги, ризики, план впровадження.
    %\item Домашнє завдання: дослідити ISO 27001 A.9 та порівняти з політикою доступу.
\end{enumerate} \\ \hline

\end{tabular}
\end{table}
% -- Кінець таблиці ---

Розробка базової політики передбачає заповнення шаблону політики керування доступом для малого офісу на 10 осіб. Визначення трьох рівнів доступу (публічний для веб-сайту та прайсів, внутрішній для робочих документів, конфіденційний для фінансів та кадрів) формує розуміння класифікації інформації.
Розподіл ролей користувачів демонструє практичну реалізацію принципу need-to-know: директор має повний доступ для управлінського контролю, бухгалтер - до фінансів та кадрів для виконання функцій, менеджери - до клієнтської бази для роботи з клієнтами, звичайні працівники - лише до власних проектів.
Встановлення правил паролів (мінімум 8 символів, зміна кожні 90 днів, 3 спроби входу) закладає основи парольної політики. Міні-кейс про першій день нового працівника моделює реальний процес onboarding, включаючи послідовність кроків надання доступу.

Комплексна система управління вводить професійний інструмент - матрицю RACI (Responsible, Accountable, Consulted, Informed) для процесу надання доступу. Розподіл ролей між директором, IT-відділом, HR та керівником відділу демонструє складність корпоративних процесів: керівник відділу відповідальний за запит (R), директор затверджує (A), HR консультує (C), IT виконує технічну частину.
Розробка процедури звільнення через чек-лист деактивації (8-10 пунктів) формує розуміння критичності своєчасного закриття доступу. Оцінка ризиків несвоєчасної деактивації через три сценарії (крадіжка даних, саботаж, репутаційні втрати) підкреслює важливість процедури offboarding.

Високий рівень моделює розробку повноцінної системи Identity Management для компанії 100+ осіб. Інтеграція з Active Directory демонструє використання корпоративних стандартів Microsoft. Впровадження Single Sign-On (SSO) показує сучасні підходи до балансу між безпекою та зручністю.
Бюджетування впровадження включає розрахунок витрат на ліцензії програмного забезпечення, навчання персоналу, аудит безпеки та річну підтримку з визначенням періоду окупності (ROI). Це формує розуміння економічних аспектів інформаційної безпеки.
Презентація для керівництва (5 слайдів) розвиває навички бізнес-комунікації: ідентифікація проблем, обґрунтування рішення, демонстрація переваг, візуалізація ризиків невпровадження, поетапний план реалізації. 



\subsubsection{Практична робота №2 ``Застосування шифрування на практиці''}

Практична робота №2 присвячена формуванню практичних навичок застосування криптографічних методів захисту інформації. Розроблена система диференційованих завдань реалізує комплексний підхід до вивчення криптографії через призму професійних інтересів учнів та забезпечує перехід від теоретичних знань до практичного застосування криптографічних інструментів.

Практична робота виконується протягом одного уроку. На початку заняття відбувається актуалізація опорних знань для встановлення необхідних інструментів та формування єдиного технологічного середовища. Учні можуть використовувати як онлайн-інструменти (cryptii.com для візуального шифрування, cyberchef.org як криптографічний комбайн, replit.com для програмування), так і локальні засоби (GPG, OpenSSL, Python з бібліотекою cryptography) залежно від технічних можливостей закладу освіти.

Академічний напрямок ``Дослідницька лабораторія'' спрямований на формування дослідницьких компетентностей у галузі криптографії через експериментальне порівняння та аналіз криптографічних алгоритмів. Завдання (табл. \ref{tab:lab2_acad_task})  побудовані за принципом наукового дослідження з документуванням результатів у лабораторному журналі.


% Визначимо колір для заголовка таблиці
\definecolor{TableHeaderBlue}{RGB}{46, 82, 141}

% -- Початок таблиці ---
% \noindent запобігає відступу, а [!] просить LaTeX ``спробувати''
% розмістити таблицю тут.
\begin{table}[h!]
\caption{Завдання для академічного напрямку ``Дослідницька лабораторія''.}
\label{tab:lab2_acad_task} % Мітка для посилань на таблицю
\noindent % Забирає відступ зліва
% Налаштування для списків всередині таблиці:
% nosep - прибирає зайві вертикальні відступи
% leftmargin=* - стандартний відступ для списку
\setlist[enumerate]{nosep, leftmargin=*}


% p{width} - колонка з фіксованою шириною та авто-переносом
% >{\raggedright\arraybackslash} - вирівнювання по лівому краю (краще для списків)
\begin{tabular}{
  | >{\raggedright\arraybackslash}p{0.15\textwidth} 
  | >{\raggedright\arraybackslash}p{0.80\textwidth} |
}
\hline

% -- Кольоровий заголовок ---
% \rowcolor - колір фону, \color - колір тексту
\rowcolor{TableHeaderBlue}
\hfil \color{white}\textbf{РІВЕНЬ} & 
\hfil \color{white}\textbf{ЗАВДАННЯ} \\ \hline

\textbf{Середній} (6 балів) &
\textbf{Базове дослідження хеш-функцій::}
\begin{enumerate}
    \item Порівняння MD5, SHA-1, SHA-256 для однакового тексту
    \item Аналіз довжини результатів та часу виконання
    \item Експеримент з лавинним ефектом (зміна 1 символу)
    \item Документування спостережень у лабораторному журналі
\end{enumerate} \\ \hline

\textbf{Достатній} (9 балів)&
\textbf{Криптоаналіз (додатково до середнього рівня)::}
\begin{enumerate}
    \item  Атака на слабкі паролі через Rainbow Tables
    \item Злом трьох MD5-хешів простих паролів
    \item Частотний аналіз шифру Цезаря
    \item Визначення зсуву та розшифровка повідомлення
\end{enumerate} \\ \hline

\textbf{Високий} (12 балів)&
\textbf{Експертний рівень (додатково до достатнього рівня):}
\begin{enumerate}
    \item Експеримент зі стеганографією (приховування даних)
    \item Аналіз зміни розміру файлу та візуальних змін
    \item Створити презентацію ``Методи шифрування: слабкості та переваги''
\end{enumerate} \\ \hline
\end{tabular}
\end{table}
% -- Кінець таблиці ---

На середньому рівні учні проводять базове порівняльне тестування хеш-функцій MD5, SHA-1 та SHA-256, документуючи результати в лабораторному журналі. Експеримент з лавинним ефектом (зміна одного символу вхідних даних призводить до кардинальної зміни хешу) формує розуміння фундаментальної властивості криптографічних хеш-функцій. Використання онлайн-інструменту CyberChef забезпечує візуалізацію процесів та доступність експерименту.

Достатній рівень передбачає проведення криптоаналізу через практичну демонстрацію вразливості слабких паролів. Використання Rainbow Tables для злому MD5-хешів простих паролів (``password'', ``123456'', ``12345678'') наочно демонструє небезпеку використання застарілих алгоритмів та тривіальних паролів. Частотний аналіз класичного шифру Цезаря розвиває навички криптоаналізу та розуміння історичної еволюції криптографії.

На високому рівні учні досліджують стеганографію як альтернативний метод приховування інформації. Експеримент включає приховування текстового повідомлення в зображенні з аналізом змін розміру файлу та можливостей виявлення прихованої інформації. Результати оформлюються як презентація наукового дослідження з метою, методологією, результатами експериментів та практичними рекомендаціями для закладу освіти.



Технічний напрямок ``Лабораторія кодування'' орієнтований на формування практичних навичок створення та використання криптографічних інструментів через програмування та роботу з професійними утилітами 



\definecolor{TableHeaderBlue}{RGB}{46, 82, 141}
\definecolor{TableBodyGreen}{RGB}{229, 248, 230}
\begin{table}[h!]
\caption{Завдання для технічного напрямку ``Лабораторія кодування''.}
\label{tab:lab2_tech_task} % Мітка для посилань на таблицю
\noindent % Забирає відступ зліва
% Налаштування для списків всередині таблиці:
% nosep - прибирає зайві вертикальні відступи
% leftmargin=* - стандартний відступ для списку
\setlist[enumerate]{nosep, leftmargin=*}

% p{width} - колонка з фіксованою шириною та авто-переносом
% >{\raggedright\arraybackslash} - вирівнювання по лівому краю (краще для списків)
\begin{tabular}{
  | >{\raggedright\arraybackslash}p{0.15\textwidth} 
  | >{\raggedright\arraybackslash}p{0.8\textwidth} |
}
\hline

% -- Кольоровий заголовок ---
% \rowcolor - колір фону, \color - колір тексту
\rowcolor{TableHeaderBlue}
\hfil\color{white}\textbf{РІВЕНЬ} & 
\hfil\color{white}\textbf{ЗАВДАННЯ} \\ \hline

% -- Рівень I ---
% \cellcolor - встановлює колір для окремої комірки
\cellcolor{TableBodyGreen}
\textbf{Середній} (6 балів) &
\textbf{GPG базова практика::}
\begin{enumerate}
    \item Симетричне шифрування файлів через GPG або CyberChef
    \item Використання AES з режимом CBC
    \item Шифрування 3 файлів різними паролями
\end{enumerate}
\\ \hline

% -- Рівень II ---
\cellcolor{TableBodyGreen}
\textbf{Достатній} (9 балів) &
\textbf{Python криптографія:}
\begin{enumerate}
    \item Реалізація функцій encrypt\_file() та decrypt\_file()
    \item Генерація ключа з пароля через SHA-256
    \item Створення простого менеджера паролів з хешуванням
\end{enumerate} \\ \hline

% -- Рівень III ---
\cellcolor{TableBodyGreen}
\textbf{Високий} (12 балів) &
\textbf{Проект ``Secure Messenger'':}
\begin{enumerate}
    \item Реалізація класу SecureMessenger
    \item Генерація RSA ключових пар
    \item Гібридне шифрування (AES для даних, RSA для ключів)
    \item Цифровий підпис та верифікація
    \item Збереження історії в зашифрованому вигляді
    
\end{enumerate} \\ \hline

\end{tabular}
\end{table}
% -- Кінець таблиці ---

Базовий рівень формує навички роботи з утилітою GPG (GNU Privacy Guard) або її онлайн-аналогами. Учні опановують симетричне шифрування файлів з використанням алгоритму AES в режимі CBC (Cipher Block Chaining). Завдання передбачає шифрування трьох різних файлів з різними паролями, що демонструє практичні аспекти управління ключами. Документування процесу формує навички технічної документації.

На достатньому рівні учні переходять до програмування криптографічних функцій мовою Python з використанням бібліотеки cryptography. Реалізація функцій шифрування та дешифрування файлів включає генерацію криптографічного ключа з пароля через хеш-функцію SHA-256. Створення простого менеджера паролів демонструє практичне застосування хешування для безпечного зберігання облікових даних. Тестування поведінки системи при введенні неправильного пароля формує розуміння важливості обробки помилок у криптографічних системах.

Високий рівень передбачає розробку прототипу захищеного месенджера з повним циклом криптографічного захисту. Проект включає реалізацію гібридного шифрування, що поєднує переваги симетричної (швидкість AES для великих обсягів даних) та асиметричної (безпечний обмін ключами через RSA) криптографії. Реалізація цифрового підпису забезпечує автентифікацію та цілісність повідомлень. Збереження історії переписки в зашифрованому вигляді демонструє принципи захисту даних у стані спокою. Опціональне створення графічного інтерфейсу через tkinter або веб-інтерфейсу розвиває навички створення користувацьких додатків.





Управлінський напрямок ``Лабораторія криптополітики'' спрямований на формування компетентностей з планування та впровадження криптографічного захисту на організаційному рівні (табл. \ref{tab:lab2_manag_task}) .

\definecolor{TableHeaderBrown}{RGB}{142, 85, 34}
\definecolor{TableBodyPeach}{RGB}{254, 237, 203}
\begin{table}[h!]
\caption{Завдання для управлінського напрямку ``Лабораторія криптополітики''.}
\label{tab:lab2_manag_task} % Мітка для посилань на таблицю
\noindent % Забирає відступ зліва
% Налаштування для списків всередині таблиці
\setlist[enumerate]{nosep, leftmargin=*}

\begin{tabular}{
  | >{\raggedright\arraybackslash}p{0.15\textwidth} 
  | >{\raggedright\arraybackslash}p{0.8\textwidth} |
}
\hline

% -- Кольоровий заголовок ---
\rowcolor{TableHeaderBrown}
\hfil \color{white}\textbf{РІВЕНЬ} & 
\hfil \color{white}\textbf{ЗАВДАННЯ} \\ \hline

% -- Рівень I ---
\cellcolor{TableBodyPeach}
\textbf{Середній} (6 балів) &
\textbf{Криптографічний аудит:}
    \begin{enumerate}
        \item Заповнення чек-листа аудиту для умовного стартапу
        \item Оцінка поточного стану: паролі, файли, комунікації, сертифікати
        \item Практична перевірка 3 сайтів на SSL Labs
        \item Визначення рівня захисту організації
    \end{enumerate}
 \\ \hline

% -- Рівень II ---
\cellcolor{TableBodyPeach}
\textbf{Достатній} (9 балів) &
\textbf{План криптографічної трансформації:}
\begin{enumerate}
    \item Аналіз поточного стану (AS-IS) з виявленням проблем
    \item Проектування цільового стану (TO-BE)
    \item Розробка 3-етапного плану міграції
    \item Ідентифікація ризиків та планування мітигації
\end{enumerate} \\ \hline

% -- Рівень III ---
\cellcolor{TableBodyPeach}
\textbf{Високий} (12 балів) &
\textbf{Кейс ``CryptoPolicy для медичної клініки'':}
\begin{enumerate}
    \item Розробка криптографічної політики (2 сторінки)
    
    \item Створення стандартних операційних процедур (SOP)
    \item Підготовка навчальних матеріалів для персоналу
    \item Визначення KPI та метрик безпеки
    \item Розрахунок бюджету впровадження
\end{enumerate} \\ \hline

\end{tabular}
\end{table}
% -- Кінець таблиці ---

На середньому рівні учні проводять криптографічний аудит умовного стартапу на 10 осіб через структурований чек-лист. Оцінюється поточний стан захисту паролів (алгоритми хешування, використання солі), критичних файлів (наявність шифрування, управління ключами), комунікацій (захист email, використання VPN) та веб-ресурсів (SSL-сертифікати). Практична перевірка реальних сайтів через сервіс SSL Labs формує навички технічного аудиту безпеки. Визначення загального рівня захисту організації розвиває здатність до комплексної оцінки стану інформаційної безпеки.

Достатній рівень передбачає розробку плану криптографічної трансформації організації. Аналіз поточного стану (AS-IS) виявляє критичні проблеми: використання застарілого MD5 для паролів, відсутність шифрування файлів та резервних копій, незахищена електронна пошта. Проектування цільового стану (TO-BE) включає вибір сучасних криптографічних рішень з урахуванням балансу безпеки та вартості. Розробка поетапного плану міграції (6 тижнів) демонструє розуміння складності організаційних змін. Ідентифікація ризиків переходу та планування заходів мітигації формує навички управління проектами в галузі інформаційної безпеки.

На високому рівні учні розробляють повний пакет документів для впровадження криптографічного захисту в медичній клініці. Криптографічна політика визначає сферу застосування, відповідальних осіб, стандарти алгоритмів (мінімальні та рекомендовані для різних призначень), життєвий цикл ключів та процедури реагування на компрометацію. Стандартні операційні процедури (SOP) деталізують покрокові інструкції для генерації, обміну та ротації ключів. Навчальні матеріали включають презентацію для медичного персоналу, інструкцію шифрування медичних карт та FAQ для пацієнтів. Визначення ключових показників ефективності (відсоток зашифрованих медкарт, кількість інцидентів, час на шифрування) та розрахунок бюджету впровадження формують комплексне розуміння управління криптографічним захистом на організаційному рівні.

Розроблена система завдань забезпечує практичне засвоєння криптографічних методів захисту інформації через експериментальну діяльність (академічний напрямок), програмування та використання професійних інструментів (технічний напрямок) та організаційне планування (управлінський напрямок). Поєднання різних форм діяльності – від лабораторних експериментів до розробки організаційних політик – формує цілісне розуміння ролі криптографії в системі інформаційної безпеки.




\subsubsection{Практична робота №3 ``Конфігурування міжмережевого екрану брандмауера''}

Практична робота №3 присвячена формуванню комплексних компетентностей у сфері мережевої безпеки через опанування принципів функціонування та конфігурування міжмережевих екранів. Розроблена система диференційованих завдань забезпечує інтеграцію теоретичних знань про периметровий захист із практичними навичками налаштування правил фільтрації трафіку, враховуючи професійні інтереси учнів та сучасні вимоги до фахівців у галузі кібербезпеки.

Практична робота виконується протягом уроку із можливістю додаткової самостійної роботи для завдань підвищеної складності. На початку уроку виконується актуалізація опорних знань: розуміння архітектури міжмережевих екранів, типів правил фільтрації та моделей безпеки периметра. Подальша робота передбачає практичне конфігурування брандмауера в симуляційному середовищі або на реальних системах із використанням iptables для Linux або Windows Firewall. Завдання достатнього та високого рівнів дозволяють учням поглибити розуміння комплексних сценаріїв захисту мережевої інфраструктури.

Академічний напрямок ``Аналітик систем безпеки'' спрямований на формування компетентностей з аналізу нормативно-правових вимог та розробки організаційних політик щодо мережевої безпеки. Завдання (табл. \ref{tab:lab3_acad_task}) цього напрямку акцентують увагу на документуванні конфігурацій, створенні політик безпеки та забезпеченні відповідності стандартам.


% Визначимо колір для заголовка таблиці
\definecolor{TableHeaderBlue}{RGB}{46, 82, 141}

% -- Початок таблиці ---
% \noindent запобігає відступу, а [!] просить LaTeX ``спробувати''
% розмістити таблицю тут.
{\small \setlength \tabcolsep {0.4\tabcolsep }
\begin{table}[h!]
\caption{Завдання для академічного напрямку ``Аналітик з інформаційної безпеки''.}
\label{tab:lab3_acad_task} % Мітка для посилань на таблицю
\noindent % Забирає відступ зліва
% Налаштування для списків всередині таблиці:
% nosep - прибирає зайві вертикальні відступи
% leftmargin=* - стандартний відступ для списку
\setlist[enumerate]{nosep, leftmargin=*}


% p{width} - колонка з фіксованою шириною та авто-переносом
% >{\raggedright\arraybackslash} - вирівнювання по лівому краю (краще для списків)
\begin{tabular}{
  | >{\raggedright\arraybackslash}p{0.15\textwidth} 
  | >{\raggedright\arraybackslash}p{0.80\textwidth} |
}
\hline

% -- Кольоровий заголовок ---
% \rowcolor - колір фону, \color - колір тексту
\rowcolor{TableHeaderBlue}
\hfil \color{white}\textbf{РІВЕНЬ} & 
\hfil \color{white}\textbf{ЗАВДАННЯ} \\ \hline

\small \textbf{Середній} 

(6 балів) &
\small \textbf{Базове дослідження:}
\begin{enumerate}
    \item  Визначення: Дати визначення поняттю «міжмережевий екран» та пояснити його призначення власними словами.
    \item Класифікація: Скласти таблицю або схему типів брандмауерів (мінімум 2 різновиди, наприклад: мережевий/периметровий та локальний/хостовий) з прикладами реалізації.
    \item Аналіз налаштувань:
    \begin{itemize}
        \item Перевірити статус firewall на власному ПК (Windows Defender або sudo ufw status у Linux).
        \item Зробити скріншот та прокоментувати: які з'єднання дозволені/блоковані за замовчуванням.
    \end{itemize}    
    
\end{enumerate} \\ \hline

\small \textbf{Достатній} (9 балів)&
\small \textbf{Поглиблений аналіз:}
\begin{enumerate}
    \item  Політика фільтрації: Розробити набір з 5 правил для захисту ПК. Оформити таблицею: опис правила (що дозволяє/блокує) + обґрунтування.
    \item Порівняння конфігурацій: Проаналізувати відмінності між стандартними налаштуваннями (наприклад, Windows vs Linux UFW або Ваші правила vs NIST). Виділити 2-3 ключові відмінності.
    \item Вплив на атаки: Обрати 2 кіберзагрози (сканування портів, DoS тощо) та описати, як брандмауер допомагає їм запобігти.
\end{enumerate} \\ \hline

\small \textbf{Високий} (12 балів)&
\small \textbf{Творче дослідження}
\begin{enumerate}
    \item Інфографіка: Створити схему «Архітектура мережевого захисту навчального закладу» (показати сегменти мережі та розміщення брандмауерів).
    \item Презентація (5-7 слайдів): Тема «Типові помилки при конфігуруванні брандмауерів». Навести 3-5 помилок, реальні кейси та рекомендації.
    \item Практичні рекомендації: Сформулювати 3 конкретні поради для школи щодо підвищення безпеки через налаштування firewall.
\end{enumerate} \\ \hline
\end{tabular}
\end{table}
% -- Кінець таблиці ---
}
На середньому рівні учні формують базове розуміння концепції міжмережевих екранів через самостійне формулювання визначень та створення класифікації. Робота з визначенням поняття ``міжмережевий екран'' розвиває навички аналізу та синтезу інформації з різних джерел. Створення класифікаційної таблиці типів брандмауерів (мережевий/периметровий на прикладі Cisco ASA або Fortinet, локальний/хостовий на прикладі Windows Firewall або iptables) формує системне бачення архітектури мережевого захисту. Аналіз поточних налаштувань брандмауера на власній системі через команди sudo ufw status verbose для Linux або перегляд консолі Windows Defender Firewall забезпечує практичне знайомство з реальними конфігураціями. Документування результатів скріншотами з коментарями формує навички технічної документації.

Достатній рівень передбачає поглиблений аналіз політик безпеки через розробку власного набору правил фільтрації. Створення політики з 5 правил для типового робочого місця (дозволити вихідний веб-трафік HTTP/HTTPS, заборонити всі вхідні нові з'єднання з Інтернету, дозволити відповіді на запити через stateful inspection, дозволити локальну мережеву взаємодію для принтерів та файлових серверів, блокувати P2P-протоколи) демонструє розуміння принципів мінімальних привілеїв та захисту в глибину. Порівняльний аналіз конфігурацій різних систем (Windows Firewall використовує підхід білого списку для вхідного трафіку, Linux UFW - більш гнучкі ланцюжки правил) або власних правил із рекомендаціями NIST розвиває критичне мислення. Дослідження механізмів протидії конкретним загрозам (як правила блокування діапазонів портів запобігають скануванню, як обмеження швидкості з'єднань протидіє DoS-атакам) формує розуміння практичної цінності брандмауерів.

На високому рівні учні створюють комплексні інформаційні продукти для представлення результатів дослідження. Розробка інфографіки архітектури мережевого захисту навчального закладу візуалізує розміщення брандмауерів між сегментами мережі (кабінет інформатики, адміністративна мережа, гостьовий Wi-Fi, серверна зона) та потоки дозволеного/забороненого трафіку між ними. Підготовка презентації про типові помилки конфігурування (залишені відкритими небезпечні порти 135-139 для NetBIOS, занадто широкі правила на кшталт ``permit any any'', відсутність логування, вимкнений firewall за замовчуванням, використання правил без опису призначення) із реальними прикладами інцидентів формує розуміння наслідків некоректної конфігурації. Розробка практичних рекомендацій для освітнього закладу (впровадження політики ``default deny'' для вхідного трафіку, сегментація мережі з окремими правилами для кожної зони, регулярний аудит правил брандмауера) демонструє здатність застосовувати теоретичні знання для вирішення практичних завдань.




Технічний напрямок ``Мережевий інженер безпеки'' орієнтований на формування практичних навичок конфігурування, налаштування та тестування міжмережевих екранів. Завдання побудовані за принципом поступового ускладнення від базових правил фільтрації до комплексних сценаріїв захисту.
\definecolor{TableHeaderBlue}{RGB}{46, 82, 141}
\definecolor{TableBodyGreen}{RGB}{229, 248, 230}

% Глобальні налаштування списків (або можна залишити локально перед таблицею)
\setlist[enumerate]{nosep, leftmargin=*}

\begin{longtable}{
  | >{\raggedright\arraybackslash}p{0.13\textwidth} 
  | >{\raggedright\arraybackslash}p{0.8\textwidth} |
}

% --- ЗАГОЛОВОК ТАБЛИЦІ (CAPTION) ---
\caption{Завдання для технічного напрямку ``Мережевий інженер безпеки''.}
\label{tab:lab3_tech_task} \\

% --- ШАПКА ПЕРШОЇ СТОРІНКИ ---
\hline
\rowcolor{TableHeaderBlue}
\hfil\color{white}\textbf{РІВЕНЬ} & 
\hfil\color{white}\textbf{ЗАВДАННЯ} \\ \hline
\endfirsthead

% --- ШАПКА НАСТУПНИХ СТОРІНОК (якщо таблиця розірветься) ---
\multicolumn{2}{r}{{Продовження табл. \thetable}} \\
\hline
\rowcolor{TableHeaderBlue}
\hfil\color{white}\textbf{РІВЕНЬ} & 
\hfil\color{white}\textbf{ЗАВДАННЯ} \\ \hline
\endhead

% --- НИЖНЯ ЧАСТИНА (FOOTER) ---
\hline
\endlastfoot

% --- ТІЛО ТАБЛИЦІ ---

% -- Рівень I ---
\rowcolor{white} % Скидання кольору рядка перед cellcolor
\cellcolor{TableBodyGreen}
\small \textbf{Середній} (6 балів) &
\cellcolor{white} % Якщо потрібно, щоб права колонка була білою, або приберіть цей рядок
\small \textbf{Базова практика:}
Симулятор або VM (Linux/Windows):
\begin{enumerate}
    \item Активувати брандмауер та встановити політику за замовчуванням (Default Deny Incoming). Команда: \texttt{sudo ufw default deny incoming}, \texttt{sudo ufw default allow outgoing}
    \item Налаштувати правило дозволу (наприклад, SSH 22/tcp або HTTP 80/tcp).
    \item Додати правило для заборони ICMP (ping) або небезпечного порту (Telnet 23).
    \item Перевірити дію правила (спробувати пінгувати або підключитися).
\end{enumerate}
\\ \hline

% -- Рівень II ---
\cellcolor{TableBodyGreen}
\small \textbf{Достатній} (9 балів) &
\small \textbf{Python криптографія / Аудит:}
\begin{enumerate}
    \item Аудит відкритих портів: \vspace{-0.5em}
    \begin{itemize}
        \item виконати сканування своєї системи з мережі, щоб побачити “до” і “після” конфігурації брандмауера. \vspace{-0.5em}
        \item запустити nmap проти основної машини, або використати онлайн-сканер портів. Які порти показані відкритими? \vspace{-0.5em}
        \item застосувати посилені правила: закрити всі непотрібні порти. Наприклад, якщо було відкрито 3–4 порти, залишити лише ті, що необхідні (22 для SSH, 80 для веб, решту – закрити).
    \end{itemize} 
    
    \item Розширити конфігурацію, використану на середньому рівні, додавши принаймні дві додаткові умови/функції: \vspace{-0.5em}
    \begin{itemize}
        \item Фільтрація за IP-адресами: створити правило, що дозволяє доступ до певного сервісу лише з конкретної IP-адреси або підмережі. \vspace{-0.5em}
        \item Логування: увімкнути логування для заблокованих пакетів.
    \end{itemize} 
    \item Створити скрипт (.sh файл) або файл конфігурації для брандмауера, який при запуску застосує всі налаштування, виконані в пунктах 1–2.
    
\end{enumerate} \\ \hline

% -- Рівень III ---
\cellcolor{TableBodyGreen}
\small \textbf{Високий} (12 балів) &
\small \textbf{Проект ``Secure Messenger'':}
\begin{enumerate}
    \item Розгорнути віртуальну мережу з декількох вузлів або використати мережевий симулятор для створення схеми “Інтернет – Firewall – Внутрішня мережа”. Налаштувати правила на Firewall, щоб: \vspace{-0.5em}
    \begin{itemize}
        \item дозволити користувачам внутрішньої мережі виходити в Інтернет (HTTP, HTTPS, DNS), але блокувати будь-які вхідні ініціативні з’єднання ззовні до внутрішніх ПК; \vspace{-0.5em}
        \item дозволити лише специфічні запити ззовні, наприклад, доступ до веб-сервера.
    \end{itemize} 
    \item Здійснити випробування налаштованої мережі: з імітованого зовнішнього вузла (інтернет) спробувати провести сканування внутрішньої мережі або перевірити доступ до закритих сервісів.
    \item Записати відео (3-5 хв) або провести живий показ роботи системи захисту.
    
\end{enumerate} \\ \hline

\end{longtable}

Базовий рівень формує фундаментальні навички роботи з брандмауером операційної системи. Увімкнення та базове налаштування UFW в Linux через команди \texttt{sudo ufw enable}, \texttt{sudo ufw default deny incoming}, \texttt{sudo ufw default allow outgoing} або активація Windows Firewall для відповідного мережевого профілю забезпечує початковий рівень захисту. Додавання правил дозволу для критичних сервісів (SSH на порту 22 для віддаленого адміністрування, HTTP/HTTPS на портах 80/443 для веб-доступу, DNS на порту 53 для розрішення доменних імен) демонструє принцип явного дозволу необхідних сервісів. Тестування конфігурації через спроби підключення зсередини та ззовні мережі формує навички верифікації налаштувань безпеки. Документування процесу у вигляді технічного звіту з командами та їх виводом розвиває навички професійної документації.

На достатньому рівні учні опановують розширені можливості брандмауерів та методи аналізу безпеки. Проведення сканування портів інструментом nmap до та після конфігурування (\texttt{nmap -sS target\_ip} для SYN-сканування) наочно демонструє зменшення поверхні атаки через закриття непотрібних портів. Реалізація фільтрації за IP-адресами (команда \texttt{sudo ufw allow from 192.168.1.0/24 to any port 22} дозволяє SSH тільки з локальної підмережі) формує розуміння контекстно-залежного контролю доступу. Увімкнення логування (\texttt{sudo ufw logging on}) з подальшим аналізом записів у \texttt{/var/log/ufw.log} розвиває навички моніторингу безпеки. Створення shell-скрипта автоматизації, що містить всі команди налаштування з коментарями, демонструє здатність до автоматизації рутинних операцій та забезпечення відтворюваності конфігурацій.

Високий рівень передбачає комплексне моделювання захисту корпоративної мережі. Створення в симуляторі Packet Tracer або через віртуальні машини схеми ``Інтернет - Firewall - DMZ - Внутрішня мережа'' із відповідними ACL (Access Control Lists) демонструє розуміння сегментації мережі. Налаштування правил, що дозволяють користувачам внутрішньої мережі виходити в Інтернет через NAT, але блокують всі вхідні ініціативні з'єднання, реалізує принцип асиметричного захисту. Розміщення веб-сервера в DMZ з дозволом доступу тільки на порти 80/443 демонструє концепцію демілітаризованої зони. Тестування сценарію ``атака та захист'' через спроби сканування портів з ``зовнішньої'' машини та перевірку блокування брандмауером формує навички пентестингу. Створення відеодемонстрації або детальної презентації з поясненням логіки правил та демонстрацією їх роботи в режимі реального часу розвиває вміння презентувати технічні рішення.

Управлінський напрямок ``Координатор мережевої безпеки'' спрямований на формування компетентностей з управління проєктами впровадження систем мережевого захисту та координації взаємодії технічних і бізнес-підрозділів (табл. \ref{tab:lab3_manag_task}).

\definecolor{TableHeaderBrown}{RGB}{142, 85, 34}
\definecolor{TableBodyPeach}{RGB}{254, 237, 203}
\begin{table}[h!]
\caption{Завдання для управлінського напрямку ``Координатор мережевої безпеки''.}
\label{tab:lab3_manag_task} % Мітка для посилань на таблицю
\noindent % Забирає відступ зліва
% Налаштування для списків всередині таблиці
\setlist[enumerate]{nosep, leftmargin=*}

\begin{tabular}{
  | >{\raggedright\arraybackslash}p{0.15\textwidth} 
  | >{\raggedright\arraybackslash}p{0.8\textwidth} |
}
\hline

% -- Кольоровий заголовок ---
\rowcolor{TableHeaderBrown}
\hfil \color{white}\textbf{РІВЕНЬ} & 
\hfil \color{white}\textbf{ЗАВДАННЯ} \\ \hline

% -- Рівень I ---
\cellcolor{TableBodyPeach}
\small \textbf{Середній} 

(6 балів) &
\small \textbf{Базове планування:}
    \begin{enumerate}
        \item Знайти вимоги щодо захисту мережі в офіційному документі (Закон про кібербезпеку / Стандарти НБУ). Законспектувати 2-3 тези.
        \item Розробити таблицю перевірки брандмауера для малої фірми (8-10 пунктів, наприклад: ``Чи закриті зайві порти?'', ``Чи ведеться логування?'').
        \item Написати чернетку ``Політики мережевого доступу'' (5-6 зрозумілих пунктів для персоналу).
    \end{enumerate}
 \\ \hline

% -- Рівень II ---
\cellcolor{TableBodyPeach}
\small \textbf{Достатній} (9 балів) &
\small \textbf{Аналіз рішень:}
\begin{enumerate}
    \item Заповнити порівняльну таблицю ``Вбудовані засоби vs Апаратний NGFW''. Критерії: Вартість, Функціональність, Складність, Масштабованість. Надати рекомендацію для малого бізнесу.
    \item Скласти алгоритм дій при інциденті (наприклад, помилка в конфігурації firewall призвела до витоку).
    \item Обрати стандарт (ISO 27001 / PCI DSS), знайти вимоги до firewall та описати дії для їх виконання.
\end{enumerate} \\ \hline

% -- Рівень III ---
\cellcolor{TableBodyPeach}
\small \textbf{Високий} (12 балів) &
\small \textbf{Стратегічне бачення:}
\begin{enumerate}
    \item Розробити архітектуру безпеки для компанії з 3 офісами (схема розміщення брандмауерів: периметр, філії, хмара).
    \item Підготувати бізнес-презентацію ``Обґрунтування інвестицій у модернізацію''. Переконати керівництво виділити бюджет (опис ризиків, ROI, вигоди).
    \item Розробити процедуру внесення змін у конфігурацію (хто подає запит, хто затверджує, як тестувати).
\end{enumerate} \\ \hline

\end{tabular}
\end{table}
% -- Кінець таблиці ---

На середньому рівні учні аналізують нормативно-правову базу мережевої безпеки. Вивчення Закону України ``Про основні засади забезпечення кібербезпеки'' або постанов НБУ щодо кіберзахисту банківських систем виявляє конкретні вимоги до використання міжмережевих екранів. Конспектування ключових тез (обов'язковість засобів фільтрації на межі корпоративної мережі та Інтернету, вимоги до документування конфігурацій, необхідність регулярного перегляду правил) формує розуміння регуляторного контексту. Розробка чек-листа аудиту брандмауера для малого підприємства з 8-10 контрольними пунктами (``Чи встановлено брандмауер на всіх точках входу в мережу?'', ``Чи задокументовано всі правила з обґрунтуванням?'', ``Чи закрито всі невикористовувані порти?'', ``Чи налаштовано логування та моніторинг?'', ``Чи проводиться регулярний перегляд правил?'') розвиває навички операційного контролю. Створення проекту ``Політики мережевого доступу'' простою мовою для персоналу організації забезпечує вміння транслювати технічні вимоги в зрозумілі інструкції для користувачів.

Достатній рівень передбачає аналіз альтернативних рішень та планування заходів реагування. Порівняння варіантів впровадження (використання вбудованих засобів Linux-сервера з iptables проти придбання апаратного NGFW від Cisco або Fortinet) за критеріями вартості (відкрите ПЗ проти ліцензій та підтримки), функціональності (базова фільтрація проти IPS/IDS, контент-фільтрації, антивірусу), складності адміністрування (необхідність Linux-експертизи проти графічного інтерфейсу) та масштабованості формує навички техніко-економічного обґрунтування. Розробка плану реагування на інцидент з неправильною конфігурацією брандмауера (негайна ізоляція скомпрометованого сегмента, виправлення правил, аналіз логів для оцінки масштабу витоку, повідомлення керівництва та клієнтів за необхідності, посилення процедур change management) демонструє розуміння процесів управління інцидентами. Аналіз вимог стандарту ISO 27001 або PCI DSS щодо брандмауерів та опис конкретних дій для забезпечення відповідності розвиває компетентності з управління комплаєнсом.

На високому рівні учні розробляють стратегічні рішення для складних організаційних структур. Проектування архітектури мережевої безпеки для розподіленої організації з трьома офісами включає визначення місць розміщення брандмауерів (периметрові в кожному офісі, внутрішні між критичними сегментами, хмарні для захисту SaaS-сервісів), їх ролей (периметрові для захисту від зовнішніх загроз, внутрішні для сегментації та контролю lateral movement, хмарні для захисту cloud workloads) та взаємодії через централізоване управління. Підготовка бізнес-презентації для керівництва з обґрунтуванням інвестицій включає кількісну оцінку ризиків поточної системи (потенційні збитки від інцидентів), аналіз переваг модернізації (зменшення ймовірності інцидентів, відповідність стандартам, покращення продуктивності), розрахунок ROI з урахуванням прямих (обладнання, ліцензії) та непрямих (навчання, простої) витрат. Розробка процедури управління змінами в конфігурації брандмауерів через формалізований процес (ініціація запиту через систему тікетів, оцінка впливу на безпеку, затвердження CAB, тестування в ізольованому середовищі, впровадження у maintenance window, верифікація та документування) демонструє розуміння ITIL-практик в контексті мережевої безпеки.

Розроблена система завдань забезпечує комплексне опанування технології міжмережевих екранів через аналітичну діяльність та створення інформаційних продуктів (академічний напрямок), практичне конфігурування та тестування систем захисту (технічний напрямок) та організаційне планування впровадження (управлінський напрямок). Поєднання різних форм діяльності – від теоретичного аналізу до практичного моделювання мережевих архітектур – формує цілісне розуміння ролі брандмауерів як ключового компонента системи мережевої безпеки сучасної організації.



\subsubsection{Завдання до підсумкового уроку}

Завершення вивчення вибіркового модуля ``Інформаційна безпека'' передбачає проведення тематичного оцінювання, форма якого є гнучкою і може проходити у будь-який спосіб. На наш погляд більш доцільним є виконання комплексних інтегрованих проєктів. Це буде сприяти і досягнення цілей профільної ліцейної освіти. Як зазначалося в аналізі, зміст модуля має значний ``розрив'' між базовим (репродуктивним) матеріалом (Основи кібергігієни) та профільними більш складними темами (Криптографія, Мережевий захист). І саме проєктна діяльність буде виконувати функцію систематизації та практичної валідації отриманих знань. Запропонуємо перелік інтегрованих проєктних тем, які забезпечують комплексне застосування знань, отриманих з усіх розділів модуля:
\begin{itemize}
    \item проєкт ``Периметр-Фортеця''
    \item проєкт ``Кібер-Ревізор''
    \item проєкт ``Крипто-Обмін''
\end{itemize}

Схарактеризуємо ці проєкти.

Проєктна робота \textbf{``Периметр-Фортеця''}  присвячена архітектурному та політичному проєктуванню комплексної системи інформаційної безпеки для малої корпоративної мережі (МКМ) на прикладі юридичної компанії. Основною метою було досягнення оптимального балансу тріади Конфіденційність, Цілісність, Доступність при мінімізації економічних витрат. Методологія реалізації базувалася на застосуванні парадигми ``Заборонити за замовчуванням'' (Deny by Default) на рівні мережевого периметру та впровадженні дискреційної моделі керування доступом (DAC) на рівні внутрішнього файлового сервера. У результаті розроблено логічну конфігурацію міжмережевого екрану, що забезпечує захист від зовнішніх DoS-атак та несанкціонованого сканування портів, а також політику безпечного віддаленого доступу через VPN-тунель (OpenVPN). Проєкт демонструє інтеграцію мережевих, системних та організаційних засобів інформаційної безпеки для нейтралізації інсайдерських та зовнішніх загроз.

Мета такого проєкту полягає у обґрунтуванні, розробці та документуванні архітектури кіберзахисту малої корпоративної мережі, що забезпечує стійкість до типових зовнішніх та внутрішніх загроз, з особливим акцентом на захисті конфіденційної інформації та безпечному віддаленому доступі.

До проєкту можна поставити таке проблемне питання - Як спроєктувати економічно ефективний, але надійний мережевий периметр для малої компанії, який забезпечить конфіденційність конфіденційних даних (тріада CIA), одночасно надаючи безпечний віддалений доступ, і чи може брандмауер повністю захистити від загроз, що виникають зсередини мережі (інсайдерські ризики)?

Робота виконується групою учнів. При розробці проєкту учні зосереджуються на різних питаннях розділів вибіркового модуля. У контексті інформаційної безпеки, першочергова увага має приділятися об'єктам захисту (Розділ І). Переходячи до системного захисту (Розділ ІІ), ключовим кроком є обґрунтування відповідної моделі керування доступом. На рівні мережевого захисту (Розділ ІІІ), основним елементом є брандмауер, який функціонує за фундаментальним принципом ``Заборонити все, що явно не дозволено'' (Deny by Default). 

Результатом проєкту  є схема мережі, документ політики безпеки ``ЮрЗахист'', файл конфігурації брандмауера .

Таким чином, інтеграція правил міжмережевого екрану, політики безпеки та керування доступом на системному рівні створює надійну ``Фортецю'' для даних компанії.

Зростання кількості цільових атак, таких як фішинг, програми-вимагачі (Ransomware) та експлойти нульового дня, вимагає переходу від фрагментарного захисту до цілісної та багаторівневої стратегії. У сучасному цифровому середовищі робоча станція кінцевого користувача є однією з найбільш критичних, але водночас найуразливіших точок ІТ-інфраструктури. Створення такої інтегрованої моделі і пропонується у проєкті \textbf{``Кібер-Ревізор''}.

Мета такого проєкту полягає у розробці та практичній реалізації цілісної стратегії захисту кінцевого користувача, що охоплює організаційні, системні та мережеві засоби. 

До проєкту можна поставити таке проблемне питання - Як забезпечити надійний та інтегрований захист конфіденційних даних на кінцевій робочій станції, враховуючи необхідність протидії гібридним загрозам (соціальна інженерія, Ransomware, мережеві атаки) за допомогою мінімального, але ефективного набору засобів?

При розробці проєкту учні зосереджуються на різних питаннях розділів вибіркового модуля. Основна увага приділяється ідентифікації загроз та аналізу їхніх джерел, включаючи розробку та обґрунтування матриці ризиків, де визначаються пріоритетні загрози (Розділ І). Розмежування прав користувачів та   застосування криптографічних методів захисту дозволить учням реалізувати системний захист. Учні також розглядають правові основи та відповідальність за порушення у сфері захисту інформації (Розділ ІІ). У сфері мережевого захисту ключовим є впровадження політики ``Deny by Default''. Ця політика на практиці реалізується через Хостовий Брандмауер. Його налаштування гарантує блокування несанкціонованого трафіку (Розділ ІІІ).

Результатом проєкту є матриця ризиків, схема політики резервного копіювання, файл конфігурації брандмауера, перелік програмних засобів, які розділені за їхнім функціональним призначенням. Запропонована багаторівнева модель захисту для кінцевої робочої станції мінімізує найкритичніші ризики та забезпечує високий рівень захисту тріади CIA.

Актуальність проєкту \textbf{``Крипто-Обмін''} зумовлена стрімким розвитком цифрових технологій та постійно зростаючою кількістю кіберзагроз. У світі, де більшість комунікацій та фінансових операцій відбувається в Інтернеті, забезпечення конфіденційності, цілісності та автентичності даних є критично важливим. Проєкт пропонує практичне вирішення проблеми захисту інформації від перехоплення та фальсифікації, які є основними мережевими загрозами, безпосередньо демонструючи, як можна запобігти несанкціонованому доступу до даних, що передаються відкритими каналами. Реалізація комплексного протоколу, що поєднує симетричне шифрування (для швидкості), асиметричне шифрування (для безпечного обміну ключами) та цифровий підпис (для автентичності), сприяє формуванню цілісного інженерного розуміння архітектури безпеки, а не лише окремих алгоритмів. Крім того, оскільки алгоритми AES та RSA є міжнародними стандартами, які використовуються в реальних системах, таких як VPN, SSL/TLS та месенджери, проєкт готує учнів до роботи з інструментами, які безпосередньо застосовуються в індустрії інформаційної безпеки.

Основною метою проєкту є розробка функціонального програмного рішення, яке наочно демонструє практичне застосування сучасних криптографічних стандартів (симетричних, асиметричних алгоритмів та цифрового підпису) для забезпечення повного циклу безпечного обміну даними.

Основним проблемним запитанням над яким будуть працювати учні це Яким чином можна ефективно інтегрувати симетричне шифрування (для швидкості та конфіденційності даних), асиметричне шифрування (для безпечного обміну ключами) та цифрові підписи (для підтвердження автентичності та цілісності) у єдиний, надійний протокол обміну, і які принципи кібергігієни є критично необхідними для безпечного управління ключами у цій системі?

Проєкт ``Крипто-Обмін'' є комплексною практичною демонстрацією ключових знань, навичок та цінностей, визначених у вибірковому модулі ``Інформаційна Безпека''. Проєкт охоплює основні поняття, забезпечуючи конфіденційність (через AES та RSA), цілісність та автентичність (через цифровий підпис та GCM), вимагає від учнів усвідомлювати відповідальність за збереження власних даних і дотримуватися правил безпечного зберігання приватних ключів (Розділ І). Учні застосовують на практиці симетричний (AES-256-GCM) та асиметричний (RSA) алгоритми, що є наріжним каменем криптографічних методів захисту. Цифровий підпис (RSA) слугує механізмом ідентифікації та аутентифікації відправника. Крім того, протокол обміну ключами, де сесійний AES-ключ шифрується публічним ключем одержувача, є практичною демонстрацією керування доступом до конфіденційної інформації, привчаючи учнів використовувати засоби адміністрування (керування ключами) для захисту (Розділ ІІ). Протокол шифрування є прямим засобом захисту, що забезпечує реагування на мережеві загрози, вразливості і атаки (зокрема, від перехоплення даних), які є основною темою цього розділу. Учні розуміють, що криптографічний захист інформації створює багатошаровий захист, доповнюючи мережевий захист, що надається брандмауерами та VPN, і гарантує захист даних навіть у разі успішного перехоплення трафіку (Розділ ІІ).

Результатом проєкту є повністю функціональний, добре прокоментований програмний інструмент, який слугує практичним прототипом безпечного криптографічного протоколу, що забезпечує конфіденційність, цілісність та автентичність даних між відправником та одержувачем шляхом інтеграції симетричного шифрування (AES), асиметричного шифрування (RSA) та цифрового підпису. Розроблений програмний застосунок є дидактичним матеріалом для вивчення криптографії та кібербезпеки. Він слугує наочною демонстрацією роботи End-to-End шифрування, дозволяючи візуально простежити етапи генерації ключів, шифрування, підписання та перевірки даних. Чітко структурований код може бути розділений на окремі модулі (RSA, AES, підпис) для поглибленого вивчення кожного криптографічного примітиву. Крім того, його можна використовувати для проведення лабораторних робіт, під час яких учні свідомо змінюють дані, підписи або ключі, щоб спостерігати, як протокол безпеки виявляє порушення цілісності чи автентичності, тим самим тестуючи його надійність. Слід також наголосити, що такий програмний застосунок можна адаптувати для реальних, некритичних застосувань. Він може слугувати базовим прототипом для вбудовування функцій захищеного обміну повідомленнями у невеликі внутрішні програми або системи, які потребують захисту даних. Реалізований механізм шифрування ключів також можна використати для захищеного зберігання та обміну конфігураційними даними або обліковими даними між сервісами. Зрештою, цей код може бути перетворений на персональну утиліту для шифрування та підписання особистих даних або архівів перед їх завантаженням у хмарне сховище, підвищуючи рівень особистої кібергігієни.


\subsection{Методичні рекомендації щодо здійснення диференційованого навчання вибіркового модуля ``Інформаційна безпека'' учнів ліцею}


\subsubsection{Загальні рекомендації щодо проведення практичних робіт}

При проведені практичних робіт бажано враховувати індивідуальні особливості учнів. На  рис.~\ref{fig:individual_track_model} запропонована модель побудови індивідуальної освітньої траєкторії кожного учня. Представлена схема відображає повний цикл диференційованого навчання, починаючи від діагностики індивідуальних особливостей учня і завершуючи оцінюванням результатів. Центральною ідеєю моделі є трикомпонентна діагностика учня, що враховує його інтереси, рівень підготовки та стиль навчання. Ці три параметри формують основу для подальшої індивідуалізації освітньої траєкторії. Система передбачає двоетапний вибір: спочатку визначається напрямок (академічний для майбутніх дослідників, технічний для практиків-інженерів, управлінський для менеджерів), а потім рівень складності завдань -- від середнього до високого. Така архітектура забезпечує гнучкість методичної системи та дає змогу адаптувати освітній процес під індивідуальні потреби кожного учня, що цілком узгоджується з сучасними принципами диференційованого навчання. 


\tikzset{
    startstop/.style={
        rectangle, 
        rounded corners=8pt,
        minimum width=4cm, 
        minimum height=1.2cm,
        text centered, 
        draw=blue!60,
        fill=blue!15,
        drop shadow,
        font=\Large\bfseries
    },
    process/.style={
        rectangle, 
        rounded corners=5pt,
        minimum width=4cm, 
        minimum height=1.2cm,
        text centered, 
        draw=teal!60,
        fill=teal!10,
        drop shadow,
        font=\large\bfseries
    },
    decision/.style={
        rectangle,
        rounded corners=3pt,
        minimum width=2.8cm, 
        minimum height=0.9cm,
        text centered,
        draw=orange!60,
        fill=orange!15,
        font=\normalsize
    },
    option/.style={
        rectangle,
        rounded corners=3pt,
        minimum width=2.8cm,
        minimum height=0.9cm,
        text centered,
        draw=purple!60,
        fill=purple!15,
        font=\normalsize
    },
    arrow/.style={
        thick,
        ->,
        >=stealth,
        draw=gray!70
    },
    wideArrow/.style={
        line width=1.5pt,
        ->,
        >=stealth,
        draw=teal!70
    }
}
\begin{figure}[h]
    \centering
    
    \begin{tikzpicture}[node distance=1cm, auto]
        
        % Діагностика
        \node (diagnostic) [process, yshift=-0.5cm] {Діагностика учня};
        
        % Три критерії діагностики
        \node (interests) [decision, below left of=diagnostic, xshift=-3.5cm, yshift=-1.2cm] {Інтереси};
        \node (level) [decision, below of=diagnostic, yshift=-1.2cm] {Рівень підготовки};
        \node (style) [decision, below right of=diagnostic, xshift=4cm, yshift=-1.2cm] {Стиль навчання};
        
        % Вибір напрямку
        \node (direction) [process, below of=level, yshift=-0.5cm] {Вибір напрямку};
        
        % Три напрямки
        \node (academic) [option, below left of=direction, xshift=-3.5cm, yshift=-1.2cm] {Академічний};
        \node (technical) [option, below of=direction, yshift=-1.2cm] {Технічний};
        \node (management) [option, below right of=direction, xshift=4cm, yshift=-1.2cm] {Управлінський};
        
        % Вибір рівня
        \node (levelChoice) [process, below of=technical, yshift=-0.5cm] {Вибір рівня};
        
        % Три рівні
        \node (basic) [option, below left of=levelChoice, xshift=-2.5cm, yshift=-1.2cm] {Середній};
        \node (advanced) [option, below of=levelChoice, yshift=-1.2cm] {Достатній};
        \node (expert) [option, below right of=levelChoice, xshift=2.5cm, yshift=-1.2cm] {Високий};
        
        % Виконання завдання
        \node (execution) [process, below of=advanced, yshift=-0.5cm] {Виконання завдання};
        
        % Оцінювання
        \node (evaluation) [process, below of=execution, yshift=-0.5cm] {Оцінювання};
        
        % Стрілки - основний потік
        \draw [wideArrow] (diagnostic) -- (interests);
        \draw [wideArrow] (diagnostic) -- (level);
        \draw [wideArrow] (diagnostic) -- (style);
        \draw [wideArrow] (interests) -- (direction);
        \draw [wideArrow] (level) -- (direction);
        \draw [wideArrow] (style) -- (direction);
        \draw [wideArrow] (direction) -- (academic);
        \draw [wideArrow] (direction) -- (technical);
        \draw [wideArrow] (direction) -- (management);
        \draw [wideArrow] (academic) -- (levelChoice);
        \draw [wideArrow] (technical) -- (levelChoice);
        \draw [wideArrow] (management) -- (levelChoice);
        \draw [wideArrow] (levelChoice) -- (basic);
        \draw [wideArrow] (levelChoice) -- (advanced);
        \draw [wideArrow] (levelChoice) -- (expert);
        \draw [wideArrow] (basic) -- (execution);
        \draw [wideArrow] (advanced) -- (execution);
        \draw [wideArrow] (expert) -- (execution);
        \draw [wideArrow] (execution) -- (evaluation);
        
    \end{tikzpicture}    
    \caption{Модель побудови індивідуальної освітньої траєкторії учня.}
    \label{fig:individual_track_model}
\end{figure}

Ефективна реалізація диференційованого підходу в навчанні інформаційної безпеки потребує створення системи різнорівневих завдань, що враховують не лише рівень підготовки учнів, але й їхні професійні інтереси та майбутню спеціалізацію. Розроблена система диференційованих завдань базується на принципах, визначених у теоретичному розділі роботи, та реалізує багатовимірну модель диференціації навчання.

У процесі розробки системи завдань було реалізовано комплексний підхід до диференціації, що базується на двох ключових вимірах: диференціація за інтересами та рівнева диференціація (табл. \ref{tab:dyferentsiatsia}).

\begin{longtable}{|p{0.25\textwidth}|p{0.3\textwidth}|p{0.35\textwidth}|}
\caption{Структура диференціації завдань вибіркового модуля ``Інформаційна безпека''.}
\label{tab:dyferentsiatsia} % Мітка для посилань на таблицю
\\
% Ви можете налаштувати ширину колонок p{...} на свій смак.
% \textwidth - це загальна ширина тексту на сторінці.
\hline
% Заголовок таблиці
\textbf{Вимір диференціації} & \textbf{Характеристика} & \textbf{Реалізація в практичних роботах } \\ \hline

% Вміст таблиці
Диференціація за інтересами& Врахування майбутньої спеціалізації учнів та варіативність тематичного наповнення  & 3 напрямки з урахуванням різних аспектів теми: академічний, технічний, управлінський \\ \hline
Рівнева диференціація & Врахування рівня підготовки та здібностей & 3 рівні складності: середній, достатній, високий \\ \hline
%Змістова диференціація & Варіативність тематичного наповнення & Різні аспекти теми залежно від профілю \\ \hline
\end{longtable}

Диференціація за інтересами реалізована через створення трьох напрямків завдань, кожен з яких корелює з певною групою майбутніх спеціальностей та професійних інтересів учнів (табл. \ref{tab:charakterystyka_napriamkiv}). Академічний напрямок ``Дослідник безпеки'' орієнтований на учнів, які планують обрати гуманітарні та соціальні спеціальності~-- право, журналістику, соціологію, міжнародні відносини. Технічний напрямок ``Етичний хакер'' розрахований на майбутніх фахівців у галузі програмування, кібербезпеки та системного адміністрування. Управлінський напрямок ``Менеджер безпеки'' спрямований на учнів, які цікавляться менеджментом, економікою, бізнесом та державним управлінням. 
Диференціація за інтересами проявляється у варіативності тематичного наповнення завдань у межах теми. Кожен напрямок акцентує увагу на специфічних аспектах теми, релевантних для майбутньої професійної діяльності учнів.

\begin{table}[h!]
\caption{Порівняльна характеристика профільних напрямків.}
\label{tab:charakterystyka_napriamkiv}
\setlength \tabcolsep {0.5\tabcolsep }
\centering

% Налаштовуємо ширину колонок. 
% p{...} - вирівнювання по верху
% m{...} - вирівнювання по середині (потрібен пакет array)
% >{\raggedright\arraybackslash} - запобігає дивному переносу слів (потрібен пакет array)
\begin{tabular}{
  | >{\raggedright\arraybackslash}p{0.18\textwidth} 
  | >{\raggedright\arraybackslash}p{0.24\textwidth} 
  | >{\raggedright\arraybackslash}p{0.24\textwidth} 
  | >{\raggedright\arraybackslash}p{0.24\textwidth} |
}
\hline
% -- Заголовок таблиці ---
\textbf{Критерій} & \textbf{Академічний напрямок} & \textbf{Технічний напрямок} & \textbf{Управлінський напрямок} \\ \hline

% -- Вміст таблиці ---
\textbf{Цільова аудиторія} & Майбутні юристи, журналісти, соціологи & Майбутні програмісти, системні адміністратори & Майбутні менеджери, економісти \\ \hline
\textbf{Основний фокус} & Аналіз, систематизація, дослідження & Практичні навички, тестування, програмування & Планування, бюджетування, управління ризиками \\ \hline
\textbf{Ключові компетентності} & Критичне мислення, аналіз джерел, створення інфопродуктів & Технічні навички, програмування, пентестинг & Стратегічне планування, комплаєнс, управління проектами \\ \hline
\textbf{Типи завдань} & Класифікації, порівняльний аналіз, інфографіка & Сканування, написання скриптів, демонстрації & Матриці ризиків, бюджети, показники ефективності \\ \hline
\textbf{Кінцевий продукт} & Аналітичні звіти, презентації, інфографіка & Технічні звіти, скрипти, відео-демонстрації & Управлінські документи, плани, положення \\ \hline
\end{tabular}
\end{table}

Рівнева диференціація забезпечує поступове ускладнення завдань від середнього до високого рівня в межах кожного профільного напрямку (табл. \ref{tab:progresiia_skladnosti}). Така структура дозволяє кожному учневі працювати в зоні найближчого розвитку, поступово нарощуючи складність виконуваних завдань відповідно до власного прогресу.

\begin{table}[h!]
\caption{Прогресія складності завдань за рівнями.}
\label{tab:progresiia_skladnosti} 
\setlength \tabcolsep {0.5\tabcolsep }
\centering
% Налаштовуємо ширину колонок. m{...} - вирівнювання по середині
\begin{tabular}{|m{0.15\textwidth}|m{0.20\textwidth}|m{0.25\textwidth}|m{0.3\textwidth}|}
\hline
% -- Заголовок таблиці ---
\textbf{Рівень} & 
\textbf{Когнітивні процеси (за таксономією Блума)} & 
\textbf{Характер діяльності} & 
\textbf{Ступінь самостійності} \\ 
\hline

% -- Базовий рівень ---
\textbf{Середній} & 
Знання, розуміння, застосування & 
Репродуктивна з елементами продуктивної & 
Виконання за зразком з детальними інструкціями \\ 
\hline

% -- Поглиблений рівень ---
\textbf{Достатній} & 
Аналіз, синтез & 
Продуктивна з елементами творчої & 
Часткова самостійність з опорними схемами \\ 
\hline

% -- Експертний рівень ---
\textbf{Високий} & 
Оцінка, створення & 
Творча, дослідницька & 
Повна самостійність, власні рішення \\ 
\hline
\end{tabular}
\end{table}




Система оцінювання диференційованих завдань (табл. \ref{tab:kryterii_otsiniuvannia})
Розроблена чотирирівнева система оцінювання враховує три критерії: повноту виконання завдань, якість аналізу та оформлення результатів.

\begin{table}[h!]
% Налаштовуємо підпис
\caption{Критерії оцінювання диференційованих завдань.}
\label{tab:kryterii_otsiniuvannia}
\setlength \tabcolsep {0.5\tabcolsep }

\centering

% Налаштовуємо ширину колонок. m{...} - вирівнювання по середині
\begin{tabular}{|m{0.08\textwidth}|m{0.12\textwidth}|m{0.22\textwidth}|m{0.25\textwidth}|m{0.25\textwidth}|}
\hline
% -- Заголовок таблиці ---
\textbf{Бали} & 
\textbf{Повнота виконання} & 
\textbf{Якість аналізу} & 
\textbf{Оформлення результатів} & 
\textbf{Додаткові критерії} \\ 
\hline

% -- Рядок 10-12 ---
\textbf{10-12} & 
100\% завдань & 
Глибокий аналіз, власні висновки, критичне мислення & 
Структуровано, візуалізація, дотримання вимог & 
Креативність, наднормові елементи \\ 
\hline

% -- Рядок 7-9 ---
\textbf{7-9} & 
70-90\% завдань & 
Достатній аналіз, є висновки, логічність & 
Структуровано, незначні недоліки & 
Виконано в термін \\ 
\hline

% -- Рядок 4-6 ---
\textbf{4-6} & 
50-70\% завдань & 
Поверхневий аналіз, описовий характер & 
Неструктуровано, без візуалізації & 
Потребує доопрацювання \\ 
\hline

% -- Рядок 1-3 ---
\textbf{1-3} & 
$<$50\% завдань & 
Відсутній аналіз, фрагментарність & 
Не відповідає вимогам & 
Несамостійне виконання \\ 
\hline
\end{tabular}
\end{table}

Високий рівень (10-12 балів) передбачає виконання всіх завдань обраного рівня, глибокий аналіз з власними висновками та структуроване оформлення з візуалізацією. Достатній рівень (7-9 балів) вимагає виконання 70-90\% завдань з достатнім аналізом та структурованим оформленням, хоча можуть бути незначні недоліки. Середній рівень (4-6 балів) характеризується виконанням 50-70\% завдань з поверхневим аналізом та неструктурованим оформленням. Початковий рівень (1-3 бали) відповідає виконанню менше 50\% завдань без аналізу та з оформленням, що не відповідає вимогам.




Представлені таблиці та схеми ілюструють комплексний підхід до організації диференційованого навчання вибіркового модуля ``Інформаційна безпека''. Система забезпечує:

\begin{itemize}
    \item 
гнучкість вибору -- учні можуть обрати напрямок та рівень складності відповідно до власних потреб;
    \item 
прогресивність навчання -- поступове ускладнення завдань забезпечує розвиток компетентностей;
    \item 
практичну спрямованість -- завдання орієнтовані на реальні професійні ситуації;
    \item 
об'єктивність оцінювання -- чіткі критерії забезпечують прозорість та справедливість оцінки;
\end{itemize}

Важливим аспектом диференційованого навчання є варіативність форм представлення результатів виконання завдань. Кожен профільний напрямок передбачає специфічні формати, що відповідають характеру майбутньої професійної діяльності та розвивають відповідні компетентності (табл.~\ref{tab:formy_rezultativ}).

Академічний напрямок акцентує увагу на аналітичних та дослідницьких форматах: письмові роботи формують навички структурованого викладу думок, візуалізація даних розвиває вміння систематизувати інформацію, а підготовка презентацій сприяє формуванню комунікативних компетентностей.

Технічний напрямок орієнтований на практичні результати: технічна документація та програмний код демонструють рівень володіння фаховими інструментами, візуальні матеріали фіксують процес виконання завдань, а використання платформ на кшталт GitHub формує навички командної роботи та версіонування.

Управлінський напрямок зосереджений на документах організаційного характеру: розробка політик та планів розвиває стратегічне мислення, створення дашбордів та матриць ризиків формує аналітичні компетентності, а інтерактивні інструменти демонструють здатність автоматизувати управлінські процеси.

\begin{table}[h!]
\caption{Форми представлення результатів диференційованих завдань}
\label{tab:formy_rezultativ}
\centering
\begin{tabular}{|m{0.25\textwidth}|m{0.2\textwidth}|m{0.2\textwidth}|m{0.25\textwidth}|}
\hline
\textbf{Напрямок} & \textbf{Текстові формати} & \textbf{Візуальні формати} & \textbf{Інтерактивні формати} \\ \hline
\textbf{Академічний} & Аналітичні звіти, есе, огляди & Інфографіка, діаграми, схеми & Презентації, веб-сайти \\ \hline
\textbf{Технічний} & Технічна документація, код & Скріншоти, блок-схеми & Відео-демонстрації, GitHub \\ \hline
\textbf{Управлінський} & Політики, плани, звіти & Дашборди, матриці ризиків & Інтерактивні калькулятори \\ \hline
\end{tabular}
\end{table}

Така варіативність форм представлення результатів дозволяє учням обирати оптимальний спосіб демонстрації набутих компетентностей відповідно до власних сильних сторін та професійних інтересів, водночас забезпечуючи об'єктивність оцінювання через чіткі критерії для кожного формату.

Запропонована модель може бути адаптована для інших тем курсу зі збереженням основних принципів диференціації.

\subsubsection{Методичні рекомендації до Практичної роботи №1 }

Практична робота №1 ``Систематизація загроз та правове поле інформаційної безпеки'' є вступною до курсу і закладає фундамент для подальшого вивчення модуля. При її проведенні рекомендується дотримуватися наступних методичних принципів.

\textbf{Організація навчального простору.} На початку заняття вчитель проводить коротке опитування для визначення початкового рівня обізнаності учнів з питань інформаційної безпеки. Це дозволяє скоригувати темп роботи та обсяг консультативної підтримки для кожної групи. Як зазначають автори \cite{Srivatanakul202225}, включення активних методів навчання до курсів з безпеки програмного забезпечення суттєво підвищує залученість студентів та якість засвоєння матеріалу.

\textbf{Формування груп.} Учні об'єднуються у групи по 3-4 особи відповідно до обраного напрямку (академічний, технічний, управлінський). Важливо забезпечити можливість зміни напрямку протягом перших двох занять, якщо учень усвідомлює, що інший профіль краще відповідає його інтересам. Дослідження авторів \cite{Nikitin20204931} підтверджують, що прикладна орієнтація завдань та міждисциплінарна інтеграція є ключовими факторами підвищення мотивації учнів до навчання.

\textbf{Робота із джерелами інформації.} Для виконання завдань середнього рівня учні використовують рекомендовані вчителем джерела: офіційні звіти CERT-UA, матеріали Державної служби спеціального зв'язку та захисту інформації України, новинні публікації про кіберінциденти. Для достатнього та високого рівнів учні самостійно здійснюють пошук додаткових джерел, що розвиває навички критичного аналізу інформації та верифікації даних.

\textbf{Диференціація консультативної підтримки.} Учні середнього рівня отримують покрокові інструкції та шаблони для заповнення. Учні достатнього рівня працюють з частково структурованими матеріалами. Учні високого рівня отримують лише загальні рекомендації, що стимулює самостійність та креативність. Такий підхід узгоджується з результатами дослідження авторів \cite{Toccafondi2024307}, які продемонстрували ефективність поєднання теоретичних занять з практичними hands-on активностями у навчанні кібербезпеки.

\textbf{Презентація результатів.} На другому уроці кожна група представляє результати своєї роботи (3-5 хвилин на групу). Решта учнів заповнюють аркуші взаємооцінювання, що сприяє розвитку критичного мислення та комунікативних навичок. Автори \cite{Broll202315990} наголошують на важливості інтерактивних scaffolded активностей для ефективного засвоєння складних концепцій.

\textbf{Оцінювання.} Рекомендується використовувати комбіноване оцінювання: 60\% -- якість виконання завдання (відповідність критеріям, повнота, коректність), 20\% -- презентація результатів, 20\% -- робота в команді та участь в обговоренні. Критерії оцінювання повідомляються учням на початку роботи.

\textbf{Рефлексія.} Наприкінці заняття учні заповнюють коротку анкету самооцінювання, що дозволяє вчителю відстежувати динаміку розвитку компетентностей та коригувати подальшу роботу. Особливу увагу слід приділити зворотному зв'язку щодо складності завдань та зрозумілості інструкцій.


\subsubsection{Методичні рекомендації до Практичної роботи №2 }

Практична робота №2 ``Застосування шифрування на практиці'' є центральною у формуванні практичних навичок криптографічного захисту інформації. Практична робота проводиться після вивчення теми ``Криптографічні методи: шифрування та гешування''. Під час такої лекції учні ознайомились з принципами та класифікацією криптографічних методів, розглянули алгоритми ``Цезаря'', ``Гамування'', AES, RSA. При проведенні практичної роботи №2 рекомендується дотримуватися наступних методичних принципів, що враховують технічну складність теми та необхідність забезпечення безпечного навчального середовища.

\textbf{Підготовка технічного середовища}. Заздалегідь вчитель  перевіряє доступність онлайн-інструментів (cryptii.com, cyberchef.org, replit.com) та за можливістю встановлює локальні засоби (GPG, Python з бібліотекою cryptography) на комп'ютери в класі. Рекомендується створити резервні варіанти доступу через мобільні пристрої учнів у разі технічних проблем. 

\textbf{Формування груп. }На початку заняття учні об'єднуються в пари або мікрогрупи по 2-3 особи в межах обраного напрямку. Парна робота є оптимальною для криптографічних експериментів, оскільки дозволяє моделювати ситуації обміну зашифрованими повідомленнями між учасниками.  

\textbf{Організація експериментальної діяльності.} Для академічного напрямку необхідно підготувати шаблони лабораторних журналів (додаток Б) з чіткою структурою для документування експериментів. Учні технічного напрямку повинні мати доступ до заздалегідь підготовлених шаблонів коду (додаток Б) з коментарями, що пояснюють ключові криптографічні операції. Для управлінського напрямку слід надати приклади реальних криптографічних політик організацій (з видаленням конфіденційної інформації) (додаток Б). 

\textbf{Диференціація технічної підтримки.} Учні середнього рівня отримують детальні покрокові інструкції з візуальними підказками (скріншоти інтерфейсів, приклади команд). При виникненні труднощів вчитель демонструє процес на проекторі. Учні достатнього рівня працюють з частковими підказками - надається загальна структура коду або алгоритм дій, але конкретна реалізація залишається за учнями. Учні високого рівня отримують лише постановку задачі та посилання на документацію криптографічних бібліотек. Такий підхід узгоджується з принципами scaffolded learning, де поступове зменшення підтримки сприяє розвитку самостійності.

\textbf{Безпека експериментів. }Критично важливо наголосити учням на використанні лише навчальних даних для шифрування. Забороняється експериментувати з реальними паролями, особистими документами чи конфіденційною інформацією. Всі експерименти з криптоаналізом (злом хешів, частотний аналіз) проводяться виключно на спеціально підготовлених навчальних прикладах. Необхідно підкреслити етичні та правові аспекти використання криптографічних інструментів. 

\textbf{Демонстрація результатів.} Кожна група готує 5-хвилинну демонстрацію своїх результатів. Для академічного напрямку - це презентація ключових висновків з експериментів з візуалізацією результатів. Технічний напрямок демонструє виконання програмного коду з поясненням алгоритмів та тестуванням на прикладах. Управлінський напрямок представляє фрагменти розроблених політик з обґрунтуванням вибору криптографічних рішень. 

\textbf{Критерії оцінювання.} Рекомендується використовувати диференційоване оцінювання: для академічного напрямку - 40\% за правильність експериментів, 30\% за повноту документування, 30\% за якість аналізу результатів; для технічного напрямку - 50\% за функціональність коду, 25\% за оптимальність алгоритмів, 25\% за документування; для управлінського напрямку - 40\% за повноту політик, 30\% за обґрунтованість рішень, 30\% за практичну застосовність.

\textbf{Міжпредметні зв'язки. }Вчителю рекомендується підкреслити зв'язок криптографії з математикою (модульна арифметика, прості числа), фізикою (квантова криптографія), історією (еволюція шифрів від давнини до сучасності), правом (законодавче регулювання криптографії). Це сприяє формуванню цілісного світогляду та розумінню міждисциплінарної природи інформаційної безпеки.

\textbf{Подальший розвиток.} Наприкінці заняття учням можна запропонувати долучитися до онлайн-платформ для поглиблення знань, таких як: CryptoHack (https://cryptohack.org/) для вирішення криптографічних задач, Coursera курси з криптографії для теоретичної підготовки (https://www.coursera.org/learn/crypto), участь у CTF-змаганнях (https://ctftime.org/event/list/upcoming) з криптографічними завданнями. Рекомендується створити шкільний криптоклуб для зацікавлених учнів з регулярними зустрічами для обговорення новин криптографії та спільного вирішення складних задач. 

\textbf{Рефлексія та самооцінювання.} Наприкінці заняття учні заповнюють форму самооцінювання з питаннями: ``Який криптографічний метод здався найскладнішим і чому?'', ``Де в реальному житті ви стикаєтеся з криптографією?'', ``Які загрози може створити слабка криптографія?''. Особливу увагу слід приділити усвідомленню учнями важливості криптографії для захисту приватності в цифрову епоху. Вчитель аналізує відповіді для коригування подальшого освітнього процесу та виявлення учнів, зацікавлених у поглибленому вивченні криптографії.


\subsubsection{Методичнi рекомендацiї до Практичної роботи №3}

Практична робота №3 ``Конфігурування міжмережевого екрану брандмауера'' є ключовою у формуванні практичних навичок мережевої безпеки та потребує ретельної підготовки технічного середовища. При її проведенні рекомендується дотримуватися наступних методичних принципів, що враховують технічну складність теми та необхідність забезпечення диференційованого підходу.

\textbf{Підготовка технічного середовища. }Заздалегідь вчитель повинен перевірити доступність та працездатність необхідного програмного забезпечення. Для академічного напрямку достатньо наявності браузера для аналізу онлайн-ресурсів та офісного пакету для створення документації. Технічний напрямок потребує встановлення віртуальних машин з Linux (рекомендовано Ubuntu 22.04 LTS з попередньо встановленим UFW) або доступу до мережевих симуляторів (Cisco Packet Tracer 8.2+, GNS3). Для управлінського напрямку необхідний доступ до нормативних документів та шаблонів організаційних політик. 

\textbf{Організація навчального простору.} На початку уроку (перші 5 хвилин) проводиться експрес-діагностика готовності учнів через питання: ``Хто знайомий з поняттям firewall?'', ``Хто бачив налаштування брандмауера Windows?'', ``Хто працював з командним рядком Linux?''. Це дозволяє вчителю оперативно скоригувати рівень початкової підтримки для різних груп. Демонстрація на проекторі базових команд UFW або інтерфейсу Windows Firewall створює спільну відправну точку для всіх учнів.

\textbf{Формування робочих груп. }Учні об'єднуються в міні-групи по 2-3 особи в межах обраного напрямку, що дозволяє організувати взаємодопомогу при виконанні технічних завдань. Для технічного напрямку рекомендується формувати пари ``досвідчений - початківець'' для забезпечення peer learning. В академічному напрямку ефективні групи по 3 особи для проведення мозкового штурму при розробці класифікацій та аналізі. Управлінський напрямок може працювати як індивідуально (для розробки політик), так і в парах (для взаємного рецензування документів). Важливо забезпечити можливість консультацій між групами різних напрямків, що моделює реальну міждисциплінарну взаємодію в ІТ-проектах.

\textbf{Диференціація стартової активності.} Перші 10 хвилин практичної роботи присвячені базовому знайомству з інструментарієм відповідно до напрямку. Академічний напрямок отримує структуровану таблицю для заповнення з підказками щодо критеріїв класифікації брандмауерів. Технічний напрямок виконує покрокову інструкцію першого запуску UFW з детальними скріншотами кожного кроку. Управлінський напрямок аналізує готовий приклад чек-листа аудиту з коментарями щодо важливості кожного пункту. Така диференціація забезпечує успішний старт для всіх учнів незалежно від початкового рівня підготовки.

\textbf{Робота з технічними інструментами.} Для технічного напрямку критично важливе поетапне освоєння команд брандмауера. Спочатку учні працюють з підготовленими вчителем віртуальними машинами з базовою конфігурацією, де можна безпечно експериментувати. Команди вводяться послідовно з обов'язковою фіксацією результатів: sudo ufw status → скріншот → sudo ufw enable → скріншот → аналіз змін. Для складніших завдань (налаштування правил за IP) надаються шаблони команд з поясненням синтаксису. При роботі з Packet Tracer учні спочатку відтворюють готову топологію мережі за зразком, а потім модифікують правила ACL. 

\textbf{Організація роботи з документацією. }Академічний та управлінський напрямки інтенсивно працюють з нормативними та технічними документами. Для ефективної організації цієї роботи рекомендується підготувати папку з ключовими документами: витяги з Закону України ``Про основні засади забезпечення кібербезпеки'' з виділеними релевантними статтями, спрощену версію вимог ISO 27001 щодо мережевої безпеки, приклади реальних політик брандмауерів (знеособлені). Учні навчаються виділяти ключову інформацію через техніку ``світлофор'': зеленим - критичні вимоги, жовтим - рекомендації, червоним - заборони. Це формує навички швидкого аналізу регуляторних документів.

\textbf{Підтримка експериментальної діяльності}. Для завдань достатнього та високого рівнів, що передбачають експерименти з конфігураціями, організовується ``пісочниця'' - ізольоване середовище для безпечного тестування. Технічний напрямок отримує доступ до окремої VLAN або набору віртуальних машин, де можна моделювати атаки без ризику для основної мережі. Академічний напрямок проводить аналіз реальних конфігурацій через онлайн-сервіси перевірки SSL-сертифікатів або використовує знеособлені дампи правил iptables для аналізу. Вчитель виступає в ролі консультанта, дозволяючи учням самостійно виявляти помилки конфігурації через тестування, що формує дослідницькі навички.

\textbf{Консультативна підтримка. }Диференціюється залежно від рівня складності та напрямку. Учні середнього рівня отримують детальні покрокові інструкції з частими контрольними точками (після кожних 2-3 команд - перевірка вчителем). Учні достатнього рівня працюють з опорними схемами та підказками щодо напрямку вирішення (``Подумайте, який порт використовує SSH'', ``Перевірте синтаксис команди allow from''). Учні високого рівня отримують лише постановку задачі та посилання на man-сторінки або офіційну документацію. Важливо заохочувати взаємодопомогу в групах: той, хто виконав завдання швидше, стає тьютором для інших.

\textbf{Документування результатів. }Кожен напрямок веде специфічну документацію виконаної роботи. Технічний напрямок створює лог-файл з усіма виконаними командами та їх результатами (можна використовувати команду script для автоматичного запису сесії терміналу). Академічний напрямок оформляє аналітичну записку з таблицями, схемами та висновками. Управлінський напрямок розробляє пакет організаційних документів. Рекомендується використовувати спільні Google Docs для можливості оперативного коментування вчителем та спільної роботи в групах.

\textbf{Тестування та верифікація.} Особлива увага приділяється перевірці працездатності створених конфігурацій. Для технічного напрямку організовується взаємне тестування: одна група намагається ``атакувати'' систему іншої через сканування портів або спроби несанкціонованого доступу. Використовуються прості інструменти: telnet для перевірки доступності портів, ping для перевірки ICMP, curl для тестування веб-доступу. Академічний напрямок верифікує свої рекомендації через моделювання сценаріїв застосування. Управлінський напрямок проводить табletop exercise - словесне програвання сценарію впровадження розроблених політик.

\textbf{Презентація результатів. }Останні 10 хвилин уроку відводяться на експрес-презентації результатів. Кожна група демонструє один ключовий результат: технічний напрямок показує працюючу конфігурацію з демонстрацією блокування небажаного трафіку, академічний - найцікавіший висновок з аналізу або візуалізацію архітектури, управлінський - ключовий пункт розробленої політики. Використовується техніка ``elevator pitch'' - 60 секунд на презентацію суті роботи. Це розвиває навички стислого представлення технічних рішень.

\textbf{Оцінювання}. Використовується комплексна система оцінювання з урахуванням специфіки кожного напрямку. Для технічного напрямку: 40\% - правильність конфігурації (працездатність правил), 30\% - повнота тестування, 30\% - якість документації. Для академічного напрямку: 40\% - глибина аналізу, 30\% - якість візуалізації та структурування, 30\% - обґрунтованість висновків. Для управлінського напрямку: 40\% - відповідність документів стандартам, 30\% - реалістичність планів, 30\% - повнота охоплення аспектів безпеки. Додатково нараховуються бонусні бали за допомогу іншим групам та креативні рішення.

\textbf{Рефлексія та підсумки.} В кінці заняття учні заповнюють електронну форму самооцінювання з питаннями: ``Що було найскладнішим?'', ``Який новий навик я освоїв?'', ``Де можна застосувати отримані знання?''. Аналіз відповідей дозволяє вчителю скоригувати подальше навчання. Особливо важливо зафіксувати типові помилки конфігурування (забули зберегти правила через ufw reload, не врахували порядок правил, заблокували власний доступ) для їх обговорення на наступному занятті. 

Домашнє завдання для охочих продовжити дослідження диференціюється за напрямками. Технічний напрямок отримує завдання встановити pfSense у віртуальній машині та порівняти можливості з UFW. Академічний напрямок досліджує еволюцію технологій міжмережевих екранів від packet filtering до NGFW. Управлінський напрямок аналізує кейс реального інциденту, пов'язаного з неправильною конфігурацією брандмауера. Результати домашньої роботи презентуються на наступному занятті, створюючи зв'язність навчального процесу.

Запропонована методика проведення практичної роботи забезпечує ефективне формування компетентностей з конфігурування міжмережевих екранів через поєднання теоретичного аналізу, практичної діяльності та організаційного планування, враховуючи індивідуальні особливості та професійні інтереси учнів.


\subsection*{Висновки до 2 розділу}

\addcontentsline{toc}{subsection}{Висновки до 2 розділу}
За результатами порівняльного аналізу програм навчання інформаційної безпеки для учнів 5-11 класів обґрунтовано змістовне наповнення вибіркового модуля. Визначено, що модуль охоплює ключові теми: систематизацію загроз, правове поле інформаційної безпеки, керування доступом, криптографічний захист та безпеку в мережі Інтернет.

Із врахуванням принципу поступового ускладнення матеріалу та  наступності стосовно базового модуля шкільного курсу інформатики розроблено календарно-тематичне планування модуля. Згідно календарного плану створено систему диференційованих завдань для трьох профільних напрямків: академічного (``Дослідник безпеки'', ``Аналітик систем безпеки''), технічного (``Етичний хакер'') та управлінського (``Менеджер безпеки''). Кожен напрямок передбачає три рівні складності: середній, достатній та високий.

Для забезпечення гнучкості у виборі напрямку і рівня складності завдань запропоновано модель побудови індивідуальної освітньої траєкторії учня, що базується на трикомпонентній діагностиці (інтереси, рівень підготовки, стиль навчання). Сформульовано методичні рекомендації щодо проведення практичних робіт, зокрема стосовно організації навчального простору, формування груп, диференціації консультативної підтримки, критеріїв оцінювання та організації рефлексії.

Отже, у другому розділі розроблено часткові методики диференційованого навчання вибіркового модуля ``Інформаційна безпека'' для учнів ліцеїв: визначені мета й зміст, розроблені засоби навчання, обґрунтовані  методи й форми навчання. 
%\begin{enumerate}
%    \item Здійснено порівняльний аналіз програм навчання інформаційної безпеки для учнів 5-11 класів та обґрунтовано змістовне наповнення вибіркового модуля. Визначено, що модуль охоплює ключові теми: систематизацію загроз, правове поле інформаційної безпеки, керування доступом, криптографічний захист та безпеку в мережі Інтернет.
    
%    \item Розроблено календарно-тематичне планування модуля, що включає 8 годин аудиторної роботи та передбачає проведення 2 практичних робіт. Планування враховує принципи поступового ускладнення матеріалу та забезпечує наступність із базовим курсом інформатики.
    
%    \item Створено систему диференційованих завдань для трьох профільних напрямків: академічного (``Дослідник безпеки'', ``Аналітик систем безпеки''), технічного (``Етичний хакер'') та управлінського (``Менеджер безпеки''). Кожен напрямок передбачає три рівні складності: середній, достатній та високий.
    
%    \item Розроблено детальні завдання для Практичної роботи №1 ``Систематизація загроз та правове поле інформаційної безпеки'' та Уроку №4 ``Керування доступом: ІАА. Дискреційна/Мандатна моделі'' з урахуванням специфіки кожного напрямку та рівня складності.
    
%    \item Запропоновано модель побудови індивідуальної освітньої траєкторії учня, що базується на трикомпонентній діагностиці (інтереси, рівень підготовки, стиль навчання) та забезпечує гнучкість у виборі напрямку і рівня складності завдань.
    
%    \item Визначено форми представлення результатів диференційованих завдань для кожного напрямку: текстові формати (аналітичні звіти, технічна документація, політики), візуальні формати (інфографіка, блок-схеми, дашборди) та інтерактивні формати (презентації, відео-демонстрації, веб-ресурси).
    
%    \item Сформульовано методичні рекомендації щодо проведення практичних робіт, включаючи організацію навчального простору, формування груп, диференціацію консультативної підтримки, критерії оцінювання та організацію рефлексії.
%\end{enumerate}

%Розроблена методична система забезпечує реалізацію принципів диференційованого навчання згідно з моделлю К.~Томлінсон та враховує сучасні тенденції у галузі навчання кібербезпеки.